% concatenated from HK_CS_memory_2026_09_02.tex
\documentclass[11pt]{amsart}
\usepackage{amsfonts,amssymb, amsthm,amsmath}
\topmargin-0.1in \textwidth6.in \textheight8.5in \oddsidemargin0in
\evensidemargin0in
\usepackage{indentfirst}
\usepackage{amssymb,mathrsfs,graphicx,subfigure, enumerate}
\usepackage{amsmath,amsfonts,amscd,amsthm,bbm}
\usepackage{extpfeil}
\usepackage{graphicx,colortbl}
\usepackage{graphicx}
\usepackage{subfigure}
\usepackage {stfloats}
\usepackage{pstricks}
\usepackage{float}
\usepackage{epsfig}
\usepackage{caption}
%\usepackage{subcaption}
\usepackage[pagewise]{lineno}
\usepackage{bm}
\graphicspath{{Figure/}}
\usepackage[english]{babel}
\usepackage{tikz}
\usepackage[title]{appendix}
\usepackage[T1]{fontenc}
%\usepackage{refcheck}%必须放在hyperref之前，并且运行时候必须注释掉下面的这两个超链接命令
\usepackage[colorlinks=true]{hyperref}
\hypersetup{urlcolor=red, citecolor=blue}

\allowdisplaybreaks[4]
%%\linenumbers

%\linespread{1.5}%ÐÞ¸ÄÐоà
%%%%%%%%%%%%%%%%%%%%%%%%%%%%%%%%%%%%%%%%%%%%%%%%%%%%%%%%%%%%% table of contents
%\usepackage{tocvsec2}
%\def\subsectionname{\quad}%
%%%%%%%%%%%%%%%%%%%%%%%%%%%%%%%%%%%%%%%%%%%%%%%%%%%%%%%%%%%%%

%\usepackage[draft]{Mytheorems}
%\usepackage{Mytheorems}




\title[Consensus and Flocking with Memory]{Exponential Consensus and Flocking in Multi-Agent Systems with Infinite Fading Memory}



\author[C. Pignotti]{Cristina Pignotti}
\address[Cristina Pignotti]{\newline Dipartimento di Ingegneria e Scienze dell'Informazione e Matematica, Universit\`{a} dell'Aquila, 67100 L'Aquila, Italy}
\email{cristina.pignotti@univaq.it}

\author[Y.-Q. Wang]{Yu-Qing Wang}
\address[Yu-Qing Wang]{\newline School of Mathematical Sciences, Dalian University of Technology, Dalian 116024, China \newline \& \newline Dipartimento di Ingegneria e Scienze dell'Informazione e Matematica, Universit\`{a} dell'Aquila, 67100 L'Aquila, Italy}
\email{yqwang202309@163.com}




\newtheorem{theorem}{Theorem}[section]
%\newtheorem{lemma}{Lemma}
\newtheorem{lemma}{Lemma}[section]
\newtheorem{corollary}{Corollary}[section]
\newtheorem{proposition}{Proposition}[section]
%\newtheorem{remark}{Remark}
\newtheorem{remark}{Remark}[section]
\newtheorem{definition}{Definition}[section]
\newtheorem{assumption}{Assumption}
\newtheorem{notation}{Notation}
\numberwithin{equation}{section}

\renewcommand{\theequation}{\thesection.\arabic{equation}}
\renewcommand{\thetheorem}{\thesection.\arabic{theorem}}
\renewcommand{\thelemma}{\thesection.\arabic{lemma}}
\newcommand{\bbz}{\mathbb Z}
\newcommand{\bbs}{\mathbb S}
\newcommand{\bbp}{\mathbb P}
\newcommand{\ddiv}{\textrm{div}}
\newcommand{\bn}{\bf n}
\newcommand{\rr}[1]{\rho_{{#1}}}
%\newcommand{\thh}{\theta}
\newcommand{\D}{\mathcal D}
\newcommand{\I}{\mathcal I}
%\newcommand{\T}{\mathcal \theta}
\newcommand{\RR}{\mathcal R}
\newcommand{\oo}{\mathcal O}
\newcommand{\f}{\mathcal F}
\newcommand{\LL}{\mathcal L}
\newcommand{\TT}{\mathcal T}
\newcommand{\x}{\bm x}
\newcommand{\y}{\bm y}
\newcommand{\z}{\bm z}
\newcommand{\w}{\bm w}
\newcommand{\vv}{\bm v}
\newcommand{\e}{\mathrm{e}}
\renewcommand{\i}{\mathrm{i}}
%\DeclareMathOperator{\eu}{e}
%\DeclareMathOperator{\ramuno}{i}
\def\charf {\mbox{{\text 1}\kern-.24em {\text l}}}
\renewcommand{\arraystretch}{1.5}


\newcommand{\EE}{\mathbb E}
\newcommand{\R}{\mathbb R}
\newcommand{\CC}{\mathbb C}
\newcommand{\s}{\mathcal S}
\newcommand{\T}{\mathbb \theta}
\newcommand{\N}{\mathbb N}
\newcommand{\G}{\mathcal G}
\newcommand{\V}{\mathcal V}
\newcommand{\C}{\mathcal C}
\newcommand{\n}{\mathcal N}
\newcommand{\mcL}{\mathcal L}
\newcommand{\mcR}{\mathcal R}
\newcommand{\mcK}{\mathcal K}
\newcommand{\mcH}{\mathcal H}
\newcommand{\mcI}{\mathcal I}
\newcommand{\E}{\mathcal E}
\newcommand{\Oh}{\mathcal O}
\newcommand{\NN}{\mathbb N}
\newcommand{\PP}{\mathcal P}
\newcommand{\kk}{\tilde{\kappa}}


\def\bx{\boldsymbol x}
\def\bp{\boldsymbol p}
\def\bv{\boldsymbol v}
\def\bvar{\boldsymbol \varphi}
\def\dist{\mathsf{dist}}
\def\diag{\mathsf{diag}}
\def\Id{\mathsf{Id}}
\def\buno{\mathbf{1}}

\abovedisplayskip=12pt plus 3pt minus 9pt
\abovedisplayshortskip=0pt plus 3pt
\belowdisplayskip=12pt plus 3pt minus 9pt
\belowdisplayshortskip=7pt plus 3pt minus 4pt

\begin{document}
\date{\today}


\subjclass[2020]{34D05, 34K25, 45J05, 45M05, 92D25.}
\keywords{Alignment models, Fading memory, Consensus, Flocking.}
%\thanks{$^*$Corresponding author: Cristina Pignotti}
\thanks{\textbf{Acknowledgment.} C. Pignotti is member of Gruppo Nazionale per l'Analisi Matematica,
	la Probabilità e le loro Applicazioni (GNAMPA) of the Istituto Nazionale di
	Alta Matematica (INdAM). The work of Y.-Q. Wang was supported by the State Scholarship Fund from the China Scholarship Council in 2025 (File No. 202506060082).}


\begin{abstract}
In this paper, we study the emergent collective dynamics of multi-agent systems driven by infinite distributed fading memory of Volterra type. We establish a unified theoretical framework covering both first-order opinion consensus dynamics and second-order velocity alignment flocking kinematics. By introducing Dafermos past-history transformations, the governing integro-differential systems are reformulated into dynamical systems on an extended product Hilbert spaces. For first-order dynamics, we prove that fading memory inherently provides a hidden dissipative mechanism, guaranteeing unconditional global exponential consensus with or without instantaneous communication forces. For second-order dynamics, we obtain unconditional exponential flocking for the pure fading memory system and we give a sufficient condition for flocking when an instantaneous interaction is also present.
In particular, this condition is always satisfied, namely the flocking occurs unconditionally, when the influence function has a divergent tail.
\end{abstract}
\maketitle \centerline{\date}


\tableofcontents

\section{Introduction}

In recent years, multi-agent systems have attracted significant attention across several scientific disciplines, including robotics, economics, control theory, mathematical biology, and social sciences \cite{Bullo2009, Camazine2001, Jackson2008, Jadbabaie2003, OlfatiSaber2007}. A fundamental aspect of multi-agent dynamics is understanding how simple self-organizing interaction rules lead to the emergence of globally coordinated collective behaviors, such as opinion consensus in human societies and flocking or schooling in biological swarms \cite{Castellano2009, Vicsek1995}. To describe continuous opinion formation and decision-making processes in social networks, Hegselmann and Krause \cite{Hegselmann2002} proposed the celebrated Hegselmann--Krause (HK) bounded confidence model:
\begin{equation}\label{eq:hk_classic}
\dot{\boldsymbol{x}}_i(t) = \frac{1}{N} \sum_{k=1}^{N} a_{ik}(t)\bigl(\boldsymbol{x}_k(t) - \boldsymbol{x}_i(t)\bigr), \qquad i \in [N] := \{1, \ldots, N\},
\end{equation}
where $\boldsymbol{x}_i(t) \in \mathbb{R}^n$ represents the continuous opinion state of the $i$-th agent at time $t$, and $a_{ik}(t) = \mathbf{1}_{\{\|\boldsymbol{x}_i(t) - \boldsymbol{x}_k(t)\| \le \epsilon\}}$ denotes the state-dependent interaction weight bounded by a confidence threshold $\epsilon > 0$. In this framework, agents adjust their opinions selectively by averaging over neighbors whose opinions lie within their confidence range \cite{Canuto2012, Jabin2014}. Several  recent extensions consider always positive interactions and more general interaction coefficients \cite{Choi2021, Haskovec2021}.


On the other hand, to model second-order swarming phenomena and velocity alignment in mobile physical or biological groups, Cucker and Smale \cite{Cucker2007} established the canonical flocking paradigm, which was subsequently generalized to second-order velocity alignment dynamics governed by abstract communication functions:
\begin{equation}\label{eq:cs_classic}
\begin{cases}
\dot{\boldsymbol{x}}_i(t) = \boldsymbol{v}_i(t), \\[4pt]
\dot{\boldsymbol{v}}_i(t) = \displaystyle\frac{1}{N} \sum_{k=1}^{N} \psi(\|\boldsymbol{x}_i(t) - \boldsymbol{x}_k(t)\|)\bigl(\boldsymbol{v}_k(t) - \boldsymbol{v}_i(t)\bigr), \qquad i \in [N],
\end{cases}
\end{equation}
where $\boldsymbol{x}_i(t) \in \mathbb{R}^n$ and $\boldsymbol{v}_i(t) \in \mathbb{R}^n$ represent the spatial position and velocity vectors of agent $i$, respectively. Here, the communication function $\psi : \mathbb{R}_+ \to \mathbb{R}_+$ is an abstract, non-increasing, strictly positive function quantifying the distance-dependent interaction strength between agents \cite{Ha2009}. Both the HK opinion framework \eqref{eq:hk_classic} and second-order velocity alignment kinematics \eqref{eq:cs_classic} have been extensively investigated regarding asymptotic consensus and cluster formation \cite{Choi2021, Jabin2014}, velocity flocking across fixed and switching topologies \cite{Dong, Fanqin}, as well as their kinetic descriptions and mean-field limits \cite{Carrillo2010, Ha2009, Ha2}.

Despite the analytical success of classical HK and CS models, standard formulations rely heavily on local-in-time interactions, assuming that agent state adjustments depend strictly on instantaneous state evaluations. In practical, social, biological, and technological networks, information transmission and cognitive processing naturally involve non-negligible time lags and cognitive inertia \cite{Chiara, Choi2021, Continelli, Continelli2023, Du2025, Haskovec2021, Pignotti2018, Cartabia}. To account for temporal delay effects, several delayed multi-agent models have been proposed in the literature:
\begin{equation*}\label{eq:hk_delay}
\dot{\boldsymbol{x}}_i(t) = \sum_{k=1}^{N} a_{ik}(t)\bigl(\boldsymbol{x}_k(t - \tau(t)) - \boldsymbol{x}_i(t)\bigr),
\end{equation*}
where $\tau(t) \ge 0$ represents a pointwise or time-varying communication delay \cite{Choi2021, Continelli2023}.

However, pointwise time delays consider information from only a single past instant $t - \tau(t)$. In practical multi-agent networks, decision-making and collective interactions naturally rely on the continuous accumulation of historical experience \cite{Jiang2025, Liu2023}. In social dynamics, the history of mutual agreement or disagreement fundamentally shapes how individuals evaluate information, giving rise to an emergent collective memory \cite{Boschi2021}. Recent studies also demonstrate that individual memory capacities quantitatively govern opinion propagation and consensus formation in complex networks \cite{Liu2026}. Beyond opinion dynamics, historical memory has recently become an essential mechanism in artificial intelligence, where agents utilize workflow memory abstracted from past action trajectories to guide complex task execution \cite{wang2025agent}.

Mathematically, hereditary memory effects spanning an infinite history are classically modeled using Volterra-type integro-differential equations in viscoelasticity, heat conduction, and continuum physics (see, e.g., \cite{Chepyzhov2006, Conti2020, Dafermos1970, Giorgi2001}). In multi-agent synchronization, memory effects have recently been incorporated into phase oscillator ensembles such as the Kuramoto model \cite{Cho2025}. However, existing synchronization frameworks restrict the memory integral to the finite interval $[0, t].$ We refer to \cite{Hask_memory} for recent results about clustering in spontaneous particle aggregation with finite memory. To the best of our knowledge, the emergent dynamics of multi-agent systems subject to infinite distributed fading memory over $(-\infty, t]$ remain largely unexplored.

Inspired by Volterra's hereditary principles \cite{Volterra1930} and fading memory theories \cite{Conti2020, Dafermos1970, Giorgi2001}, this paper establishes a unified theoretical framework for multi-agent opinion consensus and second-order velocity alignment flocking with infinite fading memory. Specifically, we investigate two fundamental classes of systems: first-order opinion dynamics and second-order velocity alignment dynamics. For first-order opinion dynamics, historical relative states are accumulated through a symmetric memory kernel $g_{ik}(s)$. We formulate the following two first-order models:
\vspace{0.2cm}

\noindent\textbf{Model I (First-Order Pure Memory):}
\begin{equation}\label{eq:model1}
\left\{
\begin{aligned}
&\dot{\boldsymbol{x}}_i(t) = \sum_{k=1}^{N} \int_{0}^{\infty} g_{ik}(s)\bigl(\boldsymbol{x}_k(t-s)-\boldsymbol{x}_i(t-s)\bigr)\,ds,~~ i \in [N],\\
&\boldsymbol{x}_i(t)=\boldsymbol{x}^{in}_i(t), \quad t\in(-\infty,0].
\end{aligned}
\right.
\end{equation}
Note that the memory integral covers the entire infinite past. Therefore, the global well-posedness of the integro-differential equations inherently necessitates the initial history to be well-defined over $(-\infty, 0]$. Here, the initial data $\x_i^{in}(t),~\forall i\in[N]$ are assumed to be bounded and continuous functions on $(-\infty, 0]$.

While Model I captures dynamics driven purely by historical accumulation, practical multi-agent interactions often involve a combination of cognitive memory and immediate environmental stimuli. To account for the coexistence of both mechanisms, we incorporate instantaneous interactions, leading to:
\vspace{0.2cm}

\noindent\textbf{Model II (First-Order Instantaneous Interaction + Memory):}
{\small
\begin{equation}\label{eq:model2}
\left\{
\begin{aligned}
&\dot{\boldsymbol{x}}_i(t) = \sum_{k=1}^{N} \psi(\|\boldsymbol{x}_i(t)-\boldsymbol{x}_k(t)\|) \bigl(\boldsymbol{x}_k(t)-\boldsymbol{x}_i(t)\bigr) + \sum_{k=1}^{N}\int_{0}^{\infty} g_{ik}(s) \bigl(\boldsymbol{x}_k(t-s)-\boldsymbol{x}_i(t-s)\bigr)\,ds,\\
&\boldsymbol{x}_i(t)=\boldsymbol{x}^{in}_i(t),~~t\in(-\infty,0],
\end{aligned}
\right.
\end{equation}
}
where $\psi : \mathbb{R}_+ \to \mathbb{R}_+$ denotes a continuous, non-increasing, strictly positive communication function. The initial history $\x_i^{in}(t),~\forall i\in[N]$ satisfies the same boundedness and continuity conditions as in Model I.

Beyond first-order opinion dynamics, we extend our framework to second-order velocity alignment kinematics ($\dot{\boldsymbol{x}}_i(t) = \boldsymbol{v}_i(t)$), where the acceleration is driven by relative velocity interactions. Similar to the first-order case, we first investigate a purely memory-driven velocity alignment model:
\vspace{0.2cm}

\noindent\textbf{Model III (Second-Order Pure Memory):}
\begin{equation}\label{eq:model3}
\left\{
\begin{aligned}
&\dot{\boldsymbol{x}}_i(t) = \boldsymbol{v}_i(t), \\
&\dot{\boldsymbol{v}}_i(t) = \displaystyle\sum_{k=1}^{N} \int_{0}^{\infty} g_{ik}(s)\bigl(\boldsymbol{v}_k(t-s)-\boldsymbol{v}_i(t-s)\bigr)\,ds,~~i \in [N].\\
\end{aligned}
\right.
\end{equation}
Because the infinite memory integral exclusively acts on the velocity components, a bounded and continuous initial history is only required for the velocities, while the spatial positions necessitate only point initial values. The initial conditions are thus specified as:
\begin{equation*}\label{initial-eq-model3}
\left\{
\begin{aligned}
&\boldsymbol{x}_{i}(0) = \boldsymbol{x}_{i}^{0},\\
&\boldsymbol{v}_i(t)=\boldsymbol{v}^{in}_i(t), \quad t\in(-\infty,0].
\end{aligned}
\right.
\end{equation*}

Finally, incorporating classic state-dependent instantaneous communication yields the full mixed-interaction flocking model:
\vspace{0.2cm}

\noindent\textbf{Model IV (Second-Order Instantaneous Interaction + Memory):}
{\small
\begin{equation}\label{eq:model4}
\begin{cases}
\dot{\boldsymbol{x}}_i(t) = \boldsymbol{v}_i(t), \\[4pt]
\dot{\boldsymbol{v}}_i(t) = \displaystyle\sum_{k=1}^{N} \psi(\|\boldsymbol{x}_i(t)-\boldsymbol{x}_k(t)\|) \bigl(\boldsymbol{v}_k(t)-\boldsymbol{v}_i(t)\bigr) + \sum_{k=1}^{N} \int_{0}^{\infty} g_{ik}(s)\bigl(\boldsymbol{v}_k(t-s)-\boldsymbol{v}_i(t-s)\bigr)\,ds.
\end{cases}
\end{equation}
}
The initial conditions for Model IV identically mirror those of Model III, requiring $\boldsymbol{x}_{i}(0) = \boldsymbol{x}_{i}^{0}$ and a bounded continuous history $\boldsymbol{v}_i(t)=\boldsymbol{v}^{in}_i(t)$ for $t\in(-\infty,0]$.

A central theoretical goal of this paper is to find a general framework ensuring   multi-agent systems driven by Volterra fading memory exhibit global exponential state consensus or velocity flocking.

The mathematical analysis of systems \eqref{eq:model1}--\eqref{eq:model4} presents major analytical challenges. Because the interactions involve memory integrals extending over $(0, \infty),$ standard Gr\"onwall-type estimates or energy methods on finite-dimensional phase spaces \cite{Ha2009, Cartabia} cannot be directly applied, as the state derivative depends explicitly on the entire historical trajectory. Furthermore, in pure memory dynamics (Models I and III) where instantaneous coupling is absent ($\psi \equiv 0$), standard kinetic energy derivatives lack state or velocity dissipation terms, making it difficult to establish asymptotic convergence. To overcome these technical obstacles, we establish a novel analytical approach based on the following key ingredients:
\begin{enumerate}
    \item [(i)] \textbf{Dafermos History Variable Transformation:} Inspired by the methodology introduced by Dafermos in viscoelasticity \cite{Dafermos1970}, we define the position and velocity past-history variables $$\boldsymbol{\eta}_i^t(s) := \int_{t-s}^t \boldsymbol{x}_i(r)\,dr, \quad \boldsymbol{\xi}_i^t(s) := \int_{t-s}^t \boldsymbol{v}_i(r)\,dr.$$
     By performing integration by parts on the Volterra memory integrals, the non-autonomous integro-differential models are transformed into autonomous dynamical systems on extended Hilbert spaces governed by linear transport operators $\partial_t \boldsymbol{\eta}_i^t(s) = \boldsymbol{x}_i(t) - \partial_s \boldsymbol{\eta}_i^t(s)$ and $\partial_t \boldsymbol{\xi}_i^t(s) = \boldsymbol{v}_i(t) - \partial_s \boldsymbol{\xi}_i^t(s)$.
    \item [(ii)] \textbf{Auxiliary Multipliers and Hidden Dissipation:} To address the lack of instantaneous state dissipation in Models I and III, we construct auxiliary multiplier functionals ($\Phi(t)$ for first-order and $\Phi_v(t)$ for second-order). By differentiating these multipliers along trajectories and utilizing Cauchy--Schwarz and Young inequalities, we extract the "hidden dissipation" generated by the kernel derivative $\mu'_{ik}(s) \le -\delta \mu_{ik}(s)$, where $\mu_{ik}(s) = -g'_{ik}(s)$. Combining these multipliers with the standard Dafermos energy, we construct strictly coercive Lyapunov functionals $\mathcal{E}(t)$ and $\mathcal{E}_v(t)$ that prove unconditional global exponential convergence.
    \item [(iii)] \textbf{Lyapunov Position-Integral Functional for Flocking Dynamics:} For Model IV, the position-dependent communication weight $\psi(\|\boldsymbol{x}_i(t)-\boldsymbol{x}_k(t)\|)$ naturally decays as the spatial distance expands, which threatens to diminish the instantaneous alignment dissipation. To overcome this critical issue, instead of employing traditional bootstrap arguments, we introduce the running maximum of the spatial diameter $Y(t) := \max\limits_{s \in [0,t]} X_{\max}(s)$ and construct a novel, non-increasing Lyapunov position-integral functional $L(t)$. This functional seamlessly couples the square root of the velocity energy $\sqrt{E_v(t)}$ with the spatial capacity integral evaluated at $Y(t)$. By calculating the Dini derivatives, we establish $D^+ L(t) \le 0$. Under an explicit integral capacity condition on the initial data, this strict monotonicity inherently confines the spatial diameter to a uniform finite bound $R_*$ for all times, which guarantees exponential velocity alignment.
\end{enumerate}

The rest of this paper is organized as follows. Section \ref{sec:pre} presents preliminaries, fundamental assumptions on memory kernels and communication functions, formal mathematical definitions of state consensus and velocity flocking, and detailed proofs of the Dafermos history transformations. Section \ref{sec:main_result} summarizes all main theoretical results regarding well-posedness, exponential consensus, and flocking. In Section \ref{sec:proof_wp}, we provide the complete proof of global well-posedness for both first-order and second-order models. Section \ref{sec:proof_1st} contains the detailed proofs of unconditional global exponential consensus for Models I and II. In Section \ref{sec:proof_2nd}, we present the proofs of conditional and unconditional exponential flocking for Models IV and III. Finally, we draw conclusions in Section \ref{sec:conclusion}.


\section{Preliminaries}\label{sec:pre}



\noindent\textbf{Notation.} In this paper, for column vectors in $\mathbb{R}^n$, we adopt the standard Euclidean norm denoted by $\|\cdot\|$. We set $[N] := \{1, \ldots, N\}$. The ensemble position vector and velocity vector are denoted by $\boldsymbol{x} := (\boldsymbol{x}_1, \ldots, \boldsymbol{x}_N) \in (\mathbb{R}^n)^N$ and $\boldsymbol{v} := (\boldsymbol{v}_1, \ldots, \boldsymbol{v}_N) \in (\mathbb{R}^n)^N$, respectively. We define the maximum spatial diameter and maximum velocity diameter as
\begin{equation*}
	X_{\max}(t) := \max_{i,k\in[N]} \|\boldsymbol{x}_i(t) - \boldsymbol{x}_k(t)\|, \qquad V_{\max}(t) := \max_{i,k\in[N]} \|\boldsymbol{v}_i(t) - \boldsymbol{v}_k(t)\|.
\end{equation*}
Moreover, we denote $\mathbb{R}_+:=[0, +\infty).$ Finally, for all $t\in(-\infty,0],$ we denote the initial data by $\x^{in}(t):=\left(\x^{in}_1(t),\x^{in}_2(t),...,\x^{in}_N(t)\right),~\boldsymbol{v}^{in}(t):=\left(\boldsymbol{v}^{in}_1(t),\boldsymbol{v}^{in}_2(t),...,\boldsymbol{v}^{in}_N(t)\right).$ Especially, for $t=0,$ we further denote
$$
\x^{0}:=\left(\x^{0}_1,\x^{0}_2,...,\x^{0}_N\right):=\left(\x^{in}_1(0),\x^{in}_2(0),...,\x^{in}_N(0)\right).
$$


Now, we present the fundamental standing assumptions on the communication weight functions and memory kernels, formulate the precise mathematical definitions of consensus and flocking, and establish the Dafermos past-history transformations.

\begin{assumption}[Communication Weights]\label{assum:psi}
{\rm
The state-dependent communication function $\psi : \mathbb{R}_+ \to \mathbb{R}_+$ is globally Lipschitz continuous, non-increasing, and strictly positive.
}
\end{assumption}

\begin{assumption}[Fading Memory Kernel]\label{assum:g}
{\rm
For each pair $i,k \in [N]$, the memory kernel $g_{ik} : \mathbb{R}_+ \to \mathbb{R}_+$ satisfies:
\begin{enumerate}
  \item [\rm (i)] \textbf{Symmetry and Positivity:} For all $s \ge 0,$
  $$
  g_{ik}(s) = g_{ki}(s) \ge 0.
  $$
  \item [\rm (ii)] \textbf{Regularity and Decay:} $g_{ik}$ is twice continuously differentiable and non-increasing, with $g_{ik}(s) \rightarrow  0$ as $s\rightarrow \infty.$ Moreover, defining for all $s\ge0,$
  \begin{equation*}\label{eq:mu_def}
  \mu_{ik}(s) := -g_{ik}'(s) \ge 0,
  \end{equation*}
  there exists a uniform constant $\delta > 0$ such that for all $s\ge0,$
  \begin{equation}\label{eq:g_decay}
      \mu_{ik}'(s) + \delta\, \mu_{ik}(s) \le 0.
  \end{equation}
\end{enumerate}
}
\end{assumption}

\begin{remark}\label{rem:kernel_implication}
{\rm
By the Fundamental Theorem of Calculus and the boundary condition $g_{ik}(s) \rightarrow  0$ as $s\rightarrow \infty$, the memory kernel $g_{ik}$ is recovered as the tail integral of $\mu_{ik}$:
\begin{equation*}\label{eq:g_as_tail_of_mu}
g_{ik}(s) = \int_s^\infty \big(-g_{ik}'(\tau)\big)\,d\tau = \int_s^\infty \mu_{ik}(\tau)\,d\tau, \qquad \forall s \ge 0.
\end{equation*}
Integrating \eqref{eq:g_decay} over $[s,\infty)$ yields $-\mu_{ik}(s) + \delta\, g_{ik}(s) \le 0$, which implies the domination property
\begin{equation*} \label{eq:mu_domination}
g_{ik}(s) \le \frac{1}{\delta}\, \mu_{ik}(s), \qquad \forall s\ge0.
\end{equation*}
Gr\"onwall's inequality applied to \eqref{eq:g_decay} guarantees $\mu_{ik}(s)\le \mu_{ik}(0)e^{-\delta s}$, and hence $g_{ik}(s)\le \frac{\mu_{ik}(0)}{\delta}e^{-\delta s}$. We define the maximum total mass parameter of the original kernels by
\begin{equation*}
\kappa_{\max} := \max_{i,k\in[N]} \int_0^\infty g_{ik}(s)\,ds < \infty,
\end{equation*}
and the mass parameters of the auxiliary tail density $\mu_{ik}$ by
$$
\lambda_{ik} := \int_0^\infty \mu_{ik}(s)\,ds = g_{ik}(0) > 0, \quad \lambda_{\min} := \min_{i,k\in[N]}\lambda_{ik},
$$
$$
\lambda_{\max} := \max_{i,k\in[N]}\lambda_{ik}, \quad \mu^0_{\max} := \max_{i,k\in[N]} \mu_{ik}(0) < \infty.
$$
}
\end{remark}

Next, we introduce the formal mathematical definitions of asymptotic consensus and flocking for the multi-agent systems under study.


\begin{definition}\label{def:consensus}
{\rm
A solution $\boldsymbol{x}(t) = (\boldsymbol{x}_1(t), \ldots, \boldsymbol{x}_N(t))$ to the first-order system \eqref{eq:model1} or \eqref{eq:model2} is said to achieve \textbf{asymptotic  consensus} if the maximum spatial diameter converges to zero asymptotically:
\begin{equation*}
\lim_{t \to \infty} X_{\max}(t) = \lim_{t \to \infty} \max_{i,k\in[N]} \|\boldsymbol{x}_i(t) - \boldsymbol{x}_k(t)\| = 0.
\end{equation*}
If there exist positive constants $C_0 > 0$ and $\gamma > 0$ such that for all $t \ge 0,$
$$
X_{\max}(t) \le C_0 e^{-\gamma t},
$$
then the system achieves \textbf{exponential consensus}.
}
\end{definition}

\begin{definition}\label{def:flocking}
{\rm
A solution $(\boldsymbol{x}(t), \boldsymbol{v}(t)) = (\boldsymbol{x}_1(t), \ldots, \boldsymbol{x}_N(t), \boldsymbol{v}_1(t), \ldots, \boldsymbol{v}_N(t))$ to the second-order system \eqref{eq:model3} or \eqref{eq:model4} is said to achieve \textbf{asymptotic flocking} if the following two conditions are satisfied simultaneously:
\begin{enumerate}
    \item [\rm (i)] \textbf{Velocity Alignment:} The maximum velocity diameter decays to zero asymptotically, i.e.,
    \begin{equation*}
    	\lim_{t \to \infty}
    V_{\max}(t) = 	\lim_{t \to \infty}\max_{i,k\in[N]} \|\boldsymbol{v}_i(t) - \boldsymbol{v}_k(t)\| =0.
    \end{equation*}
    \item [\rm (ii)] \textbf{Uniform Spatial Boundedness:} The maximum spatial diameter remains uniformly bounded for all times:
    \begin{equation*}
    \sup_{t \ge 0} X_{\max}(t) = \sup_{t \ge 0} \max_{i,k\in[N]} \|\boldsymbol{x}_i(t) - \boldsymbol{x}_k(t)\| < \infty.
    \end{equation*}
\end{enumerate}
}
\end{definition}
If the maximum velocity diameter decays to zero exponentially fast,
then we say the system achieves \textbf{exponential flocking}.

Finally, we can rewrite  the integro-differential equations via the Dafermos past-history variable reformulations (see \cite{Dafermos1970}).

\begin{lemma}[First-Order History Transformation]\label{lem:uncoupled_transform_1st}
For $i \in [N]$ and $t \ge 0$, define the position past-history variable:
\begin{equation} \label{eq:eta_uncoupled}
\boldsymbol{\eta}_i^t(s) := \int_{t-s}^{t} \boldsymbol{x}_i(r)\,dr, \qquad s \ge 0.
\end{equation}
Then $\boldsymbol{\eta}_i^t(0) = \mathbf{0}$, $\partial_s \boldsymbol{\eta}_i^t(s) = \boldsymbol{x}_i(t-s)$, and the core transport identity holds strictly:
\begin{equation}\label{eq:transport_uncoupled}
\partial_t \boldsymbol{\eta}_i^t(s) = \boldsymbol{x}_i(t) - \boldsymbol{x}_i(t-s) = \boldsymbol{x}_i(t) - \partial_s \boldsymbol{\eta}_i^t(s).
\end{equation}
Furthermore, the position memory integral transforms into:
\begin{equation*}\label{eq:ibp_uncoupled}
\int_{0}^{\infty} g_{ik}(s)\bigl(\boldsymbol{x}_k(t-s)-\boldsymbol{x}_i(t-s)\bigr)\,ds = \int_0^\infty \mu_{ik}(s) \bigl(\boldsymbol{\eta}_k^t(s) - \boldsymbol{\eta}_i^t(s)\bigr)\,ds.
\end{equation*}
\end{lemma}

\begin{proof}
Identity \eqref{eq:transport_uncoupled} follows directly by differentiating \eqref{eq:eta_uncoupled} with respect to $t$. Applying integration by parts in $s,$ using $\boldsymbol{\eta}_i^t(0)=\mathbf{0}$ and $g_{ik}(s) \rightarrow  0$ as $s\rightarrow \infty$, together with the exponential decay of $g_{ik}(s)$, we have:
{\small
\begin{align*}
\int_{0}^{\infty} g_{ik}(s)\bigl(\boldsymbol{x}_k(t-s)-\boldsymbol{x}_i(t-s)\bigr)\,ds &= \int_0^\infty g_{ik}(s)\, \partial_s \bigl(\boldsymbol{\eta}_k^t(s) - \boldsymbol{\eta}_i^t(s)\bigr)\, ds \\
&= \Big[g_{ik}(s)\big(\boldsymbol{\eta}_k^t(s)-\boldsymbol{\eta}_i^t(s)\big)\Big]_0^\infty - \int_0^\infty g_{ik}'(s)\big(\boldsymbol{\eta}_k^t(s) - \boldsymbol{\eta}_i^t(s)\big)\,ds \\
&= \mathbf{0} + \int_0^\infty \big(-g_{ik}'(s)\big)\big(\boldsymbol{\eta}_k^t(s) - \boldsymbol{\eta}_i^t(s)\big)\,ds \\
&= \int_0^\infty \mu_{ik}(s) \bigl(\boldsymbol{\eta}_k^t(s) - \boldsymbol{\eta}_i^t(s)\bigr)\, ds,
\end{align*}
}
which completes the proof.
\end{proof}

%\begin{remark}[Well-definedness of the Boundary Term]\label{rem:boundary_zero}
%    {\rm
%    In the proof of Lemma \ref{lem:uncoupled_transform_1st}, the exact vanishing of the boundary term $\lim_{s\to\infty} g_{ik}(s)\boldsymbol{\eta}_i^t(s) = \mathbf{0}$ relies fundamentally on the boundedness of the initial history $\boldsymbol{x}^{in}_i(t)$. Specifically, for any fixed $t > 0$ and $s > t$, the past-history variable integrates into the negative time axis:
%    \begin{align*}
%        \boldsymbol{\eta}_i^t(s) = \int_{t-s}^t \boldsymbol{x}_i(r)\,dr = \int_{t-s}^0 \boldsymbol{x}^{in}_i(r)\,dr + \int_0^t \boldsymbol{x}_i(r)\,dr.
%    \end{align*}
%    Since $\boldsymbol{x}^{in}_i \in L^\infty(-\infty, 0; \R^n)$, let $\|\boldsymbol{x}^{in}_i\|_\infty \le M$. On the finite interval $[0, t]$, the continuous state $\boldsymbol{x}_i(r)$ is bounded by some constant $C(t)$. Taking the norm yields:
%    \begin{equation}
%        \|\boldsymbol{\eta}_i^t(s)\| \le \int_{t-s}^0 \|\boldsymbol{x}^{in}_i(r)\|\,dr + \int_0^t \|\boldsymbol{x}_i(r)\|\,dr \le M(s-t) + t C(t) \le \tilde{C} \cdot s + K,
%    \end{equation}
%    which proves that $\boldsymbol{\eta}_i^t(s)$ exhibits at most \textbf{linear growth} with respect to $s$. Under the exponential decay assumption of the fading memory kernel, it strictly holds that $\lim_{s\to\infty} s \cdot g_{ik}(s) = 0$. Consequently, the linear growth of the Dafermos variable is totally dominated by the exponential decay of the kernel, ensuring:
%    \begin{equation}
%        \lim_{s\to\infty} \|g_{ik}(s)\boldsymbol{\eta}_i^t(s)\| \le \lim_{s\to\infty} g_{ik}(s)(\tilde{C}s + K) = 0.
%    \end{equation}
%    This rigorous evaluation highlights why a bounded initial history over $(-\infty, 0]$ is a mandatory mathematical prerequisite for infinite fading memory frameworks.
%    }
%\end{remark}

\begin{lemma}[Second-Order History Transformation]\label{lem:uncoupled_transform_2nd}
For $i \in [N]$ and $t \ge 0$, define the velocity past-history variable:
\begin{equation} \label{eq:xi_def}
\boldsymbol{\xi}_i^t(s) := \int_{t-s}^{t} \boldsymbol{v}_i(r)\,dr, \qquad s \ge 0.
\end{equation}
Then $\boldsymbol{\xi}_i^t(0) = \mathbf{0}$, $\partial_s \boldsymbol{\xi}_i^t(s) = \boldsymbol{v}_i(t-s)$, and the transport identity holds:
\begin{equation}\label{eq:transport_xi}
\partial_t \boldsymbol{\xi}_i^t(s) = \boldsymbol{v}_i(t) - \boldsymbol{v}_i(t-s) = \boldsymbol{v}_i(t) - \partial_s \boldsymbol{\xi}_i^t(s).
\end{equation}
Furthermore, the uncoupled velocity memory integral transforms into:
\begin{equation}\label{eq:ibp_velocity}
\int_{0}^{\infty} g_{ik}(s)\bigl(\boldsymbol{v}_k(t-s)-\boldsymbol{v}_i(t-s)\bigr)\,ds = \int_0^\infty \mu_{ik}(s)\bigl(\boldsymbol{\xi}_k^t(s)-\boldsymbol{\xi}_i^t(s)\bigr)\,ds.
\end{equation}
\end{lemma}

\begin{proof}
Differentiating \eqref{eq:xi_def} with respect to $t$ gives \eqref{eq:transport_xi}. Applying integration by parts in $s$ analogously to Lemma \ref{lem:uncoupled_transform_1st} gives \eqref{eq:ibp_velocity}.
\end{proof}


\section{Description of Main Results}\label{sec:main_result}
In this section, we summarize our main theoretical results concerning the global well-posedness, complete state consensus, and asymptotic flocking for the multi-agent systems \eqref{eq:model1}--\eqref{eq:model4} with infinite fading memory. Precisely, we divide our results into three parts: the global existence and uniqueness of classical solutions, the exponential consensus of first-order opinion dynamics, and the asymptotic flocking of second-order velocity alignment kinematics. The complete proofs are presented in detail in Sections \ref{sec:proof_wp}--\ref{sec:proof_2nd}.

\subsection{Global Well-Posedness}
We first address the global-in-time solvability of the Cauchy problems associated with Models I--II. For first-order systems \eqref{eq:model1} and \eqref{eq:model2}, the state is prescribed by a bounded continuous initial history $\boldsymbol{\phi} \in L^\infty(-\infty, 0]; (\mathbb{R}^n)^N)$. By establishing local contraction mappings and combining them with uniform a priori Gr\"onwall estimates, we obtain the following global well-posedness theorem.

\begin{theorem}\label{thm:wellposed}
For any given bounded continuous initial history $\boldsymbol{x}^{in} \in L^\infty(-\infty, 0]; (\mathbb{R}^n)^N)$ for first-order models, there exists a unique global classical solution $\boldsymbol{x} \in C^1([0,\infty); (\mathbb{R}^n)^N)$ to Models I and II.
\end{theorem}

Analogously, for second-order systems \eqref{eq:model3} and \eqref{eq:model4}, the initial data consist of an initial position history $\boldsymbol{x}^{in} \in L^\infty(-\infty, 0]; (\mathbb{R}^n)^N)$ and an initial velocity history \\$\boldsymbol{v}^{in} \in L^\infty(-\infty, 0]; (\mathbb{R}^n)^N)$. And the global well-posedness theorem can be obtained as below.

\begin{theorem}\label{thm:wellposed-1}
For any given bounded continuous initial history \\$(\boldsymbol{x}^{0}, \boldsymbol{v}^{in}) \in L^\infty(-\infty, 0]; (\mathbb{R}^n)^N \times (\mathbb{R}^n)^N)$ for second-order models, there exists a unique global classical solution $(\boldsymbol{x},\boldsymbol{v}) \in C^1([0,\infty); (\mathbb{R}^n)^N \times (\mathbb{R}^n)^N)$ to Models III and IV.
\end{theorem}

\subsection{Consensus for First-Order Systems}
In this subsection, we present the asymptotic consensus results for the first-order opinion dynamics \eqref{eq:model1} and \eqref{eq:model2}. We first state the exponential consensus result for Model II \eqref{eq:model2}, where instantaneous communication forces and distributed fading memory act simultaneously.

\begin{theorem}[Unconditional Consensus for Model II]\label{thm:model2}
Let $\boldsymbol{x}(t)$ be the global classical solution to Model II \eqref{eq:model2}. Then $\boldsymbol{x}(t)$ achieves global exponential consensus unconditionally, i.e., there exist positive constants $C_0$ and $\gamma$ such that
\begin{equation*}
X_{\max}(t) \le C_0 e^{-\gamma t}, \qquad \forall t \ge 0.
\end{equation*}
\end{theorem}

When the instantaneous communication force is absent ($\psi \equiv 0$), Model II \eqref{eq:model2} degenerates into Model I \eqref{eq:model1}. In this regime, system \eqref{eq:model1} lacks instantaneous state dissipation. By constructing an auxiliary position multiplier functional $\Phi(t)$ to extract the hidden dissipation from the fading memory kernel, we establish that memory alone is sufficient to guarantee exponential consensus as the following theorem.

\begin{theorem}[Unconditional Consensus for Model I]\label{thm:model1}
Let $\boldsymbol{x}(t)$ be the global classical solution to Model I \eqref{eq:model1}. Then $\boldsymbol{x}(t)$ achieves global exponential consensus unconditionally, demonstrating that fading memory intrinsically induces sufficient dissipative mechanisms without instantaneous interactions.
\end{theorem}

%\begin{remark}\label{rem:consensus_discussion}
%{\rm
%Theorems \ref{thm:model2} and \ref{thm:model1} demonstrate that distributed fading memory inherently induces an effective dissipative mechanism. Unlike classical memoryless Hegselmann--Krause models where consensus requires communication graph connectivity or specific initial configurations, the fading memory acts as an intrinsic stabilizer that drives the entire ensemble to consensus unconditionally across the entire state space $(\mathbb{R}^{n})^{N}$.
%}
%\end{remark}
%
%\begin{remark}\label{rem:consensus_discussion}
%Theorems \ref{thm:model2} and \ref{thm:model1} demonstrate that distributed fading memory inherently acts as an effective dissipative mechanism. Unlike classical memoryless Hegselmann--Krause models where consensus requires communication graph connectivity or sufficiently concentrated initial opinions to avoid clustering and opinion fragmentation \cite{Blondel2009, Canuto2012, Hegselmann2002, Jabin2014, Lorenz2007}, the infinite fading memory here provides a continuous global attractive pull, driving the entire ensemble to consensus unconditionally across the entire state space $\mathbb{R}^{nN}$.
%\end{remark}

\begin{remark}\label{rem:consensus_discussion}
{\rm
Theorems \ref{thm:model2} and \ref{thm:model1} reveal that distributed fading memory acts as an intrinsic globally stabilizing dissipative mechanism in multi-agent systems.

 In the absence of memory, establishing consensus in state-dependent or switching multi-agent networks fundamentally relies on persistent connectivity conditions, such as scrambling graph structures, uniform spanning trees, or persistent chain connections over time intervals \cite{Blondel2009, Canuto2012, Su2017}. In Model I, even when instantaneous communication is entirely absent ($\psi \equiv 0$), the kernel derivative condition $\mu'_{ik}(s) \le -\delta \mu_{ik}(s)$ supplies the necessary hidden dissipation to drive the whole network to consensus without requiring any time-dependent topological connectivity assumptions.
}
\end{remark}

\subsection{Flocking for Second-Order Systems}
In this subsection, we present the asymptotic flocking results for the second-order velocity alignment systems \eqref{eq:model3} and \eqref{eq:model4}. We first define the velocity energy functional as below:
\begin{equation*}\label{eq:Ev_recall}
E_v(t) := \frac12\sum_{i=1}^N\|\boldsymbol{v}_i(t)\|^2 + \frac14\int_0^\infty\sum_{i,k=1}^N \mu_{ik}(s)\|\boldsymbol{\xi}_i^t(s)-\boldsymbol{\xi}_k^t(s)\|^2\,ds.
\end{equation*}

For Model IV \eqref{eq:model4}, where instantaneous spatial interaction $\psi$ is coupled with velocity memory, the spatial diameter $X_{\max}(t)$ may expand, leading to a decay in the interaction strength $\psi(X_{\max}(t))$. By constructing a novel Lyapunov position-integral functional $L(t)$ that couples the velocity energy with the running maximum of spatial expansion, we establish the following flocking theorem.
\begin{theorem}[Unconditional and Conditional Flocking for Model IV]\label{thm:model4}
Let $(\boldsymbol{x}(t),\boldsymbol{v}(t))$ be the global classical solution to Model IV \eqref{eq:model4}. Suppose the initial velocity energy $E_v(0)$ and the initial maximum spatial diameter $X_{\max}(0)$ satisfy the following integral capacity condition:
\begin{equation}\label{eq:flocking_capacity_cond}
\sqrt{E_v(0)} < \frac{N}{2\sqrt{2}} \int_{X_{\max}(0)}^{\infty} \min \left\{ \psi(r),\, \frac{\delta}{2N} \right\} dr.
\end{equation}
Then, the system achieves asymptotic exponential flocking.
%Specifically, the maximum spatial diameter remains uniformly bounded for all time:
%\begin{equation*}
%\sup_{t \ge 0} X_{\max}(t) \le R_* < \infty,
%\end{equation*}
%where $R_* > X_{\max}(0)$ is the unique finite root to the algebraic equation
%\begin{equation}\label{eq:R_star_def}
%\frac{N}{2\sqrt{2}} \int_{X_{\max}(0)}^{R_*} \min \left\{ \psi(r),\, \frac{\delta}{2N} \right\} dr = \sqrt{E_v(0)}.
%\end{equation}
%Furthermore, the velocity alignment decays exponentially fast:
%\begin{equation*}
%V_{\max}(t) \le 2\sqrt{2E_v(0)}\, e^{-\frac{1}{2}\gamma_v(R_*) t}, \qquad \forall t \ge 0,
%\end{equation*}
%where the uniform decay rate is given by $\gamma_v(R_*) := \min\bigl\{2N\psi(R_*),\, \delta\bigr\} > 0$.
\end{theorem}

For Model III \eqref{eq:model3}, where instantaneous spatial coupling is absent ($\psi \equiv 0$), the interaction is purely velocity-driven. By constructing an auxiliary velocity multiplier functional $\Phi_v(t)$, we prove that the velocity alignment decays exponentially fast without any constraint on the initial data, and the spatial trajectory remains uniformly bounded globally.

\begin{theorem}[Unconditional Flocking for Model III]\label{thm:model3}
Let $(\boldsymbol{x}(t),\boldsymbol{v}(t))$ be the global classical solution to Model III \eqref{eq:model3}. Then, without any restriction on the network size $N$ or kernel mass, the system globally achieves exponential velocity alignment, and the maximum spatial diameter remains uniformly bounded:
$$
\sup_{t\ge0} X_{\max}(t) < \infty.
$$
\end{theorem}

\begin{remark}\label{rem:flocking_discussion}
{\rm
We highlight several structural features of the flocking framework established in Theorems \ref{thm:model4} and \ref{thm:model3}:
\begin{enumerate}
    \item [(i)] \textbf{Unconditional vs. Conditional Dichotomy:} In Theorem \ref{thm:model4}, condition \eqref{eq:flocking_capacity_cond} provides a sharp threshold that unifies both heavy-tailed and short-range communication rates:
    \begin{itemize}
        \item If the communication function $\psi$ is non-integrable at infinity (e.g., $\psi(r) = \frac{K}{(1+r^2)^\beta}$ with $\beta \le 1/2$), the integral $\int_{X_{\max}(0)}^\infty \min\{\psi(r), \frac{\delta}{2N}\} dr$ diverges to $+\infty$. Consequently, condition \eqref{eq:flocking_capacity_cond} holds for arbitrary initial data $(X_{\max}(0), E_v(0))$, recovering unconditional global flocking.
        \item If $\psi$ is integrable at infinity (short-range interactions with $\beta > 1/2$), condition \eqref{eq:flocking_capacity_cond} requires the initial velocity fluctuation $\sqrt{E_v(0)}$ to be strictly bounded by the total spatial communication capacity, establishing conditional flocking.
    \end{itemize}
    This dichotomy perfectly matches the celebrated analytical results established for the classical non-delayed Cucker--Smale model \cite{Cucker2007, Ha2009}.
    \item [(ii)] \textbf{Role of Memory Dissipation in Model III:} In Model III, even in the complete absence of position-dependent communication forces, the fading memory alone provides sufficient friction-like dissipation in velocity space, ensuring that $\boldsymbol{v}_i(t)$ converges to the average velocity exponentially fast and preventing the spatial configuration from dispersing to infinity.
\end{enumerate}
}
\end{remark}

\section{Global Well-Posedness for both first and second order systems}\label{sec:proof_wp}
In this section, we provide the proof of Theorem \ref{thm:wellposed} and Theorem \ref{thm:wellposed-1} regarding the global-in-time existence and uniqueness of classical solutions. We present the analysis by considering the first-order opinion dynamics and the second-order flocking kinematics separately. Our general strategy relies on rewriting the integro-differential systems into equivalent Volterra-type integral equations. We first establish local well-posedness on a short time interval via the Banach Fixed-Point Theorem. Subsequently, we derive uniform a priori estimates on the state variables using Gr\"onwall's inequality to prevent finite-time blow-up, thereby extending the local solutions globally.

\noindent {\bf Proof of Theorem \ref{thm:wellposed}.} Setting
$$
X^{in} := \max_{i\in[N]} \sup_{s \le 0} \|\boldsymbol{x}^{in}_i(s)\| < \infty,
$$
we rewrite Model II \eqref{eq:model2} for $t > 0$ by splitting the memory integral over $s\in[0,t]$ and $s\in[t,\infty)$. Substituting $\tau = t-s$ on $[0,t]$ and $u = s-t$ on $[0,\infty)$ gives:
\begin{equation}\label{eq:wp_split_1st}
\begin{aligned}
\dot{\boldsymbol{x}}_i(t) &= \sum_{k=1}^N \psi(\|\boldsymbol{x}_i(t)-\boldsymbol{x}_k(t)\|) (\boldsymbol{x}_k(t)-\boldsymbol{x}_i(t)) + \sum_{k=1}^N \int_0^t g_{ik}(t-\tau)(\boldsymbol{x}_k(\tau)-\boldsymbol{x}_i(\tau)) \, d\tau \\
&\quad + \sum_{k=1}^N \int_0^\infty g_{ik}(t+u)(\boldsymbol{x}^{in}_k(-u)-\boldsymbol{x}^{in}_i(-u)) \, du.
\end{aligned}
\end{equation}
The third term is a deterministic forcing term $\boldsymbol{H}_{ik}(t) := \int_0^\infty g_{ik}(t+u)(\boldsymbol{x}^{in}_k(-u)-\boldsymbol{x}^{in}_i(-u)) \, du$, whose norm is bounded by
\begin{equation}\label{eq:H_bound_1st}
\|\boldsymbol{H}_{ik}(t)\| \le 2X^{in} \int_0^\infty g_{ik}(t+u)\,du \le 2X^{in} \kappa_{ik} \le C_H := 2X^{in} \kappa_{\max} < \infty.
\end{equation}

Integrating \eqref{eq:wp_split_1st} over $[0,t]$, we define the operator $\mathcal{L}$ acting on $\boldsymbol{x}\in C([0,T];(\mathbb{R}^n)^N)$ by
\begin{equation}\label{eq:operatorL_1st}
\begin{aligned}
[\mathcal{L}(\boldsymbol{x})]_i(t) &= \boldsymbol{x}^{in}_i(0) + \int_0^t \sum_{k=1}^N \psi(\|\boldsymbol{x}_i(r)-\boldsymbol{x}_k(r)\|)(\boldsymbol{x}_k(r)-\boldsymbol{x}_i(r)) \, dr \\
&\quad + \int_0^t \sum_{k=1}^N \int_0^r g_{ik}(r-\tau)(\boldsymbol{x}_k(\tau)-\boldsymbol{x}_i(\tau)) \, d\tau\, dr + \int_0^t \sum_{k=1}^N \boldsymbol{H}_{ik}(r) \, dr.
\end{aligned}
\end{equation}

Define $F(\boldsymbol{x}_i,\boldsymbol{x}_k):=\psi(\|\boldsymbol{x}_i-\boldsymbol{x}_k\|)(\boldsymbol{x}_k-\boldsymbol{x}_i)$. For any state vectors with
$$
\|\boldsymbol{x}_i\|,\|\boldsymbol{x}_k\|,\|\tilde{\boldsymbol{x}}_i\|,\|\tilde{\boldsymbol{x}}_k\|\le R,
$$
using the global Lipschitz continuity of $\psi$ with constant $[\psi]_{\mathrm{Lip}}$, we obtain
\begin{equation*}\label{eq:FA_lip}
\|F(\boldsymbol{x}_i,\boldsymbol{x}_k)-F(\tilde{\boldsymbol{x}}_i,\tilde{\boldsymbol{x}}_k)\| \le L_R\big(\|\boldsymbol{x}_i-\tilde{\boldsymbol{x}}_i\|+\|\boldsymbol{x}_k-\tilde{\boldsymbol{x}}_k\|\big),
\end{equation*}
where $L_R := 2R [\psi]_{\mathrm{Lip}} + \psi(0)$.

Define the Banach space $\mathcal{X}_T = C([0,T];(\mathbb{R}^n)^N)$ equipped with
$$
\|\boldsymbol{x}\|_{\mathcal{X}_T} := \sup_{t\in[0,T]}\max_{i\in[N]}\|\boldsymbol{x}_i(t)\|.
$$
Let $R := 2X^{in}$ and $B_R := \{\boldsymbol{x}\in\mathcal{X}_T: \|\boldsymbol{x}\|_{\mathcal{X}_T} \le R\}$.
For any $\boldsymbol{x} \in B_R$, we have $\|\boldsymbol{x}_k(\tau)-\boldsymbol{x}_i(\tau)\| \le 2R$. From \eqref{eq:operatorL_1st} and \eqref{eq:H_bound_1st},
\begin{align*}
\|[\mathcal{L}(\boldsymbol{x})]_i(t)\| &\le X^{in} + \int_0^t \sum_{k=1}^N \psi(0)(2R)\,dr + \int_0^t \sum_{k=1}^N \kappa_{\max}(2R)\,dr + \int_0^t \sum_{k=1}^N C_H\,dr \\
&\le X^{in} + t N \Big[ 2R \big(\psi(0)+\kappa_{\max}\big) + C_H \Big] =: X^{in} + t K_1(R).
\end{align*}
Choosing $T_1 := \frac{X^{in}}{K_1(R)}$ ensures $\mathcal{L}(B_R) \subset B_R$ for $T \le T_1$. Furthermore, for $\boldsymbol{x},\tilde{\boldsymbol{x}} \in B_R$,
{\small
\begin{align*}
\|[\mathcal{L}(\boldsymbol{x})]_i(t) - [\mathcal{L}(\tilde{\boldsymbol{x}})]_i(t)\| &\le \int_0^t \sum_{k=1}^N L_R\big(\|\boldsymbol{x}_i(r)-\tilde{\boldsymbol{x}}_i(r)\|+\|\boldsymbol{x}_k(r)-\tilde{\boldsymbol{x}}_k(r)\|\big)\,dr \\
&\quad + \int_0^t \sum_{k=1}^N \int_0^r g_{ik}(r-\tau) \big(\|\boldsymbol{x}_i(\tau)-\tilde{\boldsymbol{x}}_i(\tau)\|+\|\boldsymbol{x}_k(\tau)-\tilde{\boldsymbol{x}}_k(\tau)\|\big)\,d\tau\,dr \\
&\le T \cdot 2 N (L_R + \kappa_{\max}) \|\boldsymbol{x} - \tilde{\boldsymbol{x}}\|_{\mathcal{X}_T}.
\end{align*}
}
Setting $T_2 := \frac{1}{4 N (L_R + \kappa_{\max})}$, the operator $\mathcal{L}$ is a strict contraction on $B_R$ for $T := \min\{T_1, T_2\}$. By the Banach Fixed-Point Theorem, a unique local solution exists on $[0,T]$.

To extend the solution globally, define
$$
M(t) := \sup_{r\in[0,t]}\max_{i\in[N]}\|\boldsymbol{x}_i(r)\|.
$$
Along the local solution,
\begin{equation*}
M(t) \le X^{in} + N C_H t + \int_0^t 2 N \big(\psi(0)+\kappa_{\max}\big) M(\tau)\,d\tau.
\end{equation*}
By Gr\"onwall's inequality, for all $t \ge 0,$
\begin{equation*}\label{eq:gronwall_bound_1st}
M(t) \le \big(X^{in} + N C_H t\big) e^{2 N (\psi(0)+\kappa_{\max}) t} < \infty.
\end{equation*}
Since $M(t)$ cannot experience finite-time blow-up, standard continuation arguments confirm that $T_{\max} = \infty$. Model I \eqref{eq:model1} follows directly as a special case by setting $\psi \equiv 0$.

\noindent {\bf Proof of Theorem \ref{thm:wellposed-1}.} For Model IV \eqref{eq:model4}, let
$$
X^0 := \max_{i\in[N]}  \|\boldsymbol{x}^{0}_i\| \quad \text{and} \quad V^{in} := \max_{i\in[N]} \sup_{s \le 0} \|\boldsymbol{v}^{in}_i(s)\| < \infty.
$$
For $t > 0$, splitting the memory integral over $s\in[0,t]$ and $s\in[t,\infty)$ yields:
%\begin{equation}\label{eq:wp_split_2nd}
%\begin{aligned}
%\dot{\boldsymbol{x}}_i(t) &= \boldsymbol{v}_i(t), \\
%\dot{\boldsymbol{v}}_i(t) &= \sum_{k=1}^N \psi(\|\boldsymbol{x}_i(t)-\boldsymbol{x}_k(t)\|) (\boldsymbol{v}_k(t)-\boldsymbol{v}_i(t)) + \sum_{k=1}^N \int_0^t g_{ik}(t-\tau)(\boldsymbol{v}_k(\tau)-\boldsymbol{v}_i(\tau)) \, d\tau + \sum_{k=1}^N \boldsymbol{G}_{ik}(t),
%\end{aligned}
%\end{equation}
{\footnotesize
\begin{equation}\label{eq:wp_split_2nd}
\left\{
\begin{aligned}
\dot{\boldsymbol{x}}_i(t) &= \boldsymbol{v}_i(t),  \\
\dot{\boldsymbol{v}}_i(t) &= \sum_{k=1}^N \psi(\|\boldsymbol{x}_i(t)-\boldsymbol{x}_k(t)\|) (\boldsymbol{v}_k(t)-\boldsymbol{v}_i(t)) + \sum_{k=1}^N \int_0^t g_{ik}(t-\tau)(\boldsymbol{v}_k(\tau)-\boldsymbol{v}_i(\tau)) \, d\tau + \sum_{k=1}^N \boldsymbol{G}_{ik}(t),
\end{aligned}
\right.
\end{equation}
}
where $\boldsymbol{G}_{ik}(t) := \int_0^\infty g_{ik}(t+u)(\boldsymbol{v}^{in}_k(-u)-\boldsymbol{v}^{in}_i(-u)) \, du$ satisfies
$$
\|\boldsymbol{G}_{ik}(t)\| \le 2V^{in} \kappa_{\max} =: C_G < \infty.
$$

Integrating \eqref{eq:wp_split_2nd} over $[0,t]$, we define the solution operator $\mathcal{T}(\boldsymbol{x},\boldsymbol{v}) := (\mathcal{T}_x(\boldsymbol{x},\boldsymbol{v}), \mathcal{T}_v(\boldsymbol{x},\boldsymbol{v}))$ on $C([0,T];(\mathbb{R}^n)^N \times (\mathbb{R}^n)^N)$:
\begin{align*}
[\mathcal{T}_x(\boldsymbol{x},\boldsymbol{v})]_i(t) &= \boldsymbol{x}^0_i + \int_0^t \boldsymbol{v}_i(r)\,dr, \\
[\mathcal{T}_v(\boldsymbol{x},\boldsymbol{v})]_i(t) &= \boldsymbol{v}^{in}_i(0) + \int_0^t \sum_{k=1}^N \psi(\|\boldsymbol{x}_i(r)-\boldsymbol{x}_k(r)\|)(\boldsymbol{v}_k(r)-\boldsymbol{v}_i(r)) \, dr \\
&\quad + \int_0^t \sum_{k=1}^N \int_0^r g_{ik}(r-\tau)(\boldsymbol{v}_k(\tau)-\boldsymbol{v}_i(\tau)) \, d\tau\, dr + \int_0^t \sum_{k=1}^N \boldsymbol{G}_{ik}(r) \, dr.
\end{align*}
Define the product metric space $\mathcal{Y}_T = C([0,T];(\mathbb{R}^n)^N \times (\mathbb{R}^n)^N)$ with norm
$$
\|(\boldsymbol{x},\boldsymbol{v})\|_{\mathcal{Y}_T} := \sup_{t\in[0,T]}\max_{i\in[N]} \big( \|\boldsymbol{x}_i(t)\| + \|\boldsymbol{v}_i(t)\| \big).
$$
Setting $R := 2(X^0 + V^{in})$ and $B_R := \{(\boldsymbol{x},\boldsymbol{v}) \in \mathcal{Y}_T : \|(\boldsymbol{x},\boldsymbol{v})\|_{\mathcal{Y}_T} \le R\}$, the local Lipschitz estimates on $\psi$ and integral bounds guarantee that $\mathcal{T}$ is a strict self-mapping and contraction on $B_R$ for a sufficiently small local time horizon $T > 0$.

For global continuation, defining
$$
M_v(t) := \sup_{r\in[0,t]}\max_{i\in[N]}\|\boldsymbol{v}_i(r)\|,
$$
we derive
\begin{equation*}
M_v(t) \le V^{in} + N C_G t + \int_0^t 2 N \big(\psi(0)+\kappa_{\max}\big) M_v(\tau)\,d\tau.
\end{equation*}
By Gr\"onwall's inequality,  for all $t \ge 0,$
$$
M_v(t) \le \big(V^{in} + N C_G t\big) e^{2 N (\psi(0)+\kappa_{\max}) t} < \infty.
$$
Consequently,
$$
M_x(t) := \sup_{r\in[0,t]}\max_{i\in[N]}\|\boldsymbol{x}_i(r)\| \le X^0 + \int_0^t M_v(\tau)\,d\tau < \infty.
$$
This prevents finite-time blow-up and proves that $T_{\max} = \infty$. Model III \eqref{eq:model3} follows as a special case with $\psi \equiv 0$.


\section{Exponential Consensus for First-Order Systems}\label{sec:proof_1st}

This section is devoted to the asymptotic consensus analysis of the first-order opinion dynamics, completing the proofs of Theorem \ref{thm:model2} and Theorem \ref{thm:model1}. We begin by analyzing Model II, where the coexistence of instantaneous communication and fading memory naturally provides sufficient state dissipation, which can be directly extracted using the standard Dafermos energy functional. In contrast, Model I lacks instantaneous damping. To uncover the hidden dissipation induced exclusively by the fading memory kernel, we introduce an auxiliary cross-multiplier functional. By combining this multiplier with the standard energy, we construct a  Lyapunov functional to prove unconditional exponential consensus.

\subsection{Exponential Consensus for system \eqref{eq:model2}}
By Lemma \ref{lem:uncoupled_transform_1st}, Model II \eqref{eq:model2} maps into the autonomous extended system:
\begin{equation}\label{eq:sys_auto_model2}
\dot{\boldsymbol{x}}_i(t) = \sum_{k=1}^N \psi(\|\boldsymbol{x}_i(t)-\boldsymbol{x}_k(t)\|)\big(\boldsymbol{x}_k(t) - \boldsymbol{x}_i(t)\big) + \sum_{k=1}^N\int_0^\infty \mu_{ik}(s)\big(\boldsymbol{\eta}_k^t(s) - \boldsymbol{\eta}_i^t(s)\big)\,ds.
\end{equation}
We define the first-order position energy functional $E_0(t)$:
\begin{equation*}\label{eq:E0_def}
E_0(t) := \frac{1}{2}\sum_{i=1}^N\|\boldsymbol{x}_i(t)\|^2 + \frac{1}{4}\int_0^\infty\sum_{i,k=1}^N \mu_{ik}(s)\|\boldsymbol{\eta}_i^t(s)-\boldsymbol{\eta}_k^t(s)\|^2\,ds.
\end{equation*}

\noindent {\bf Proof of Theorem \ref{thm:model2}.} Summing \eqref{eq:sys_auto_model2} over $i \in [N]$ gives $\frac{d}{dt}\sum\limits_{i=1}^N \boldsymbol{x}_i(t) = \mathbf{0}$. By Galilean invariance, we assume without loss of generality that $\sum\limits_{i=1}^N \boldsymbol{x}_i(0) = \mathbf{0}$, whence $\sum\limits_{i=1}^N \boldsymbol{x}_i(t) \equiv \mathbf{0}$ for all $t \ge 0$.

Differentiating the state component along \eqref{eq:sys_auto_model2} and symmetrically exchanging indices $i \leftrightarrow k$ yields:
\begin{align}\label{eq:E0_dot_state}
\frac{d}{dt}\left( \frac{1}{2}\sum_{i=1}^N\|\boldsymbol{x}_i(t)\|^2 \right) &= \sum_{i=1}^N \boldsymbol{x}_i(t)^T \dot{\boldsymbol{x}}_i(t) \notag \\
&= \sum_{i=1}^N \boldsymbol{x}_i(t)^T \Biggl( \sum_{k=1}^N \psi(\|\boldsymbol{x}_i(t)-\boldsymbol{x}_k(t)\|)\big(\boldsymbol{x}_k(t) - \boldsymbol{x}_i(t)\big) \notag \\
&\qquad\qquad\qquad + \sum_{k=1}^N\int_0^\infty \mu_{ik}(s)\big(\boldsymbol{\eta}_k^t(s) - \boldsymbol{\eta}_i^t(s)\big)\,ds \Biggr) \notag \\
&= -\frac{1}{2}\sum_{i,k=1}^N \psi(\|\boldsymbol{x}_i(t)-\boldsymbol{x}_k(t)\|)\|\boldsymbol{x}_i(t)-\boldsymbol{x}_k(t)\|^2 \notag \\
&\quad -\frac{1}{2}\int_0^\infty\sum_{i,k=1}^N \mu_{ik}(s)(\boldsymbol{x}_i(t)-\boldsymbol{x}_k(t))^T(\boldsymbol{\eta}_i^t(s)-\boldsymbol{\eta}_k^t(s))\,ds.
\end{align}

Differentiating the memory energy component utilizing the transport identity \eqref{eq:transport_uncoupled}:
{\footnotesize
\begin{align}\label{eq:E0_dot_memory}
&\frac{d}{dt} \left( \frac{1}{4}\int_0^\infty\sum_{i,k=1}^N \mu_{ik}(s)\|\boldsymbol{\eta}_i^t(s)-\boldsymbol{\eta}_k^t(s)\|^2\,ds \right) \notag \\
&= \frac{1}{2}\int_0^\infty\sum_{i,k=1}^N\mu_{ik}(s)(\boldsymbol{\eta}_i^t(s)-\boldsymbol{\eta}_k^t(s))^T\bigl( \partial_t \boldsymbol{\eta}_i^t(s) - \partial_t \boldsymbol{\eta}_k^t(s) \bigr)\,ds \notag \\
&= \frac{1}{2}\int_0^\infty\sum_{i,k=1}^N\mu_{ik}(s)(\boldsymbol{\eta}_i^t(s)-\boldsymbol{\eta}_k^t(s))^T\bigl[ (\boldsymbol{x}_i(t)-\boldsymbol{x}_k(t)) - \partial_s (\boldsymbol{\eta}_i^t(s)-\boldsymbol{\eta}_k^t(s)) \bigr]\,ds \notag \\
&= \frac{1}{2}\int_0^\infty\sum_{i,k=1}^N\mu_{ik}(s)(\boldsymbol{x}_i(t)-\boldsymbol{x}_k(t))^T(\boldsymbol{\eta}_i^t(s)-\boldsymbol{\eta}_k^t(s))\,ds - \frac{1}{4}\int_0^\infty\sum_{i,k=1}^N\mu_{ik}(s)\partial_s\|\boldsymbol{\eta}_i^t(s)-\boldsymbol{\eta}_k^t(s)\|^2\,ds.
\end{align}
}

Summing \eqref{eq:E0_dot_state} and \eqref{eq:E0_dot_memory}, the cross-terms cancel identically. Applying integration by parts to the remaining $\partial_s$ term with boundary conditions $\boldsymbol{\eta}_i^t(0)=\mathbf{0}$ and $\lim\limits_{t\rightarrow\infty}\mu_{ik}(t)=0$:
{\small
\begin{align}\label{eq:dot_E0_final}
\dot{E}_0(t) &= -\frac{1}{2}\sum_{i,k=1}^N \psi(\|\boldsymbol{x}_i(t)-\boldsymbol{x}_k(t)\|)\|\boldsymbol{x}_i(t)-\boldsymbol{x}_k(t)\|^2 \notag \\
&\quad - \frac{1}{4}\sum_{i,k=1}^N \left[ \mu_{ik}(s) \|\boldsymbol{\eta}_i^t(s) - \boldsymbol{\eta}_k^t(s)\|^2 \right]_0^\infty + \frac{1}{4}\int_0^\infty\sum_{i,k=1}^N \mu'_{ik}(s)\|\boldsymbol{\eta}_i^t(s)-\boldsymbol{\eta}_k^t(s)\|^2\,ds \notag \\
&= -\frac{1}{2}\sum_{i,k=1}^N \psi(\|\boldsymbol{x}_i(t)-\boldsymbol{x}_k(t)\|)\|\boldsymbol{x}_i(t)-\boldsymbol{x}_k(t)\|^2 + \frac{1}{4}\int_0^\infty\sum_{i,k=1}^N \mu'_{ik}(s)\|\boldsymbol{\eta}_i^t(s)-\boldsymbol{\eta}_k^t(s)\|^2\,ds \le 0.
\end{align}
}

Since $E_0(t) \le E_0(0)$ for all $t \ge 0$, the maximum spatial diameter is uniformly bounded:
\begin{equation*}
X_{\max}(t) \le 2 \max_{i\in[N]} \|\boldsymbol{x}_i(t)\| \le 2\sqrt{2E_0(0)} =: R_{\max}, \qquad \forall t \ge 0.
\end{equation*}
By Assumption \ref{assum:psi}, $\psi(\|\boldsymbol{x}_i(t)-\boldsymbol{x}_k(t)\|) \ge \psi(R_{\max}) > 0$. Using $\mu'_{ik}(s) \le -\delta \mu_{ik}(s)$ from Assumption \ref{assum:g} together with the algebraic identity $\sum\limits_{i,k=1}^N \|\boldsymbol{x}_i(t)-\boldsymbol{x}_k(t)\|^2 = 2N \sum\limits_{i=1}^N \|\boldsymbol{x}_i(t)\|^2$, we obtain
\begin{align*}
\dot{E}_0(t) &\le -\psi(R_{\max}) N \sum_{i=1}^N\|\boldsymbol{x}_i(t)\|^2 - \frac{\delta}{4}\int_0^\infty\sum_{i,k=1}^N \mu_{ik}(s)\|\boldsymbol{\eta}_i^t(s)-\boldsymbol{\eta}_k^t(s)\|^2\,ds \le -\gamma E_0(t),
\end{align*}
where $\gamma := \min\{2N\psi(R_{\max}), \delta\} > 0$. Applying Gr\"onwall's Lemma gives $E_0(t) \le E_0(0)e^{-\gamma t}$, which proves unconditional exponential consensus for Model II.




\subsection{Exponential Consensus for system \eqref{eq:model1}} For Model I ($\psi \equiv 0$), equation \eqref{eq:dot_E0_final} reduces to
$$
\dot{E}_0(t) = \frac{1}{4}\int_0^\infty\sum_{i,k=1}^N \mu'_{ik}(s)\|\boldsymbol{\eta}_i^t(s)-\boldsymbol{\eta}_k^t(s)\|^2\,ds \le 0,
$$
which lacks instantaneous state dissipation. We introduce the auxiliary multiplier functional:
\begin{equation}\label{eq:multiplier_M1}
\Phi(t) := -\frac{1}{2}\int_0^\infty\sum_{i,k=1}^N\mu_{ik}(s)\big(\boldsymbol{x}_i(t)-\boldsymbol{x}_k(t)\big)^T\big(\boldsymbol{\eta}_i^t(s)-\boldsymbol{\eta}_k^t(s)\big)\,ds.
\end{equation}

\noindent {\bf Proof of Theorem \ref{thm:model1}.}
Differentiating \eqref{eq:multiplier_M1} along \eqref{eq:sys_auto_model2} (with $\psi \equiv 0$) using $\partial_t \boldsymbol{\eta}_i^t(s) = \boldsymbol{x}_i(t) - \partial_s \boldsymbol{\eta}_i^t(s)$:
\begin{align}\label{eq:Phi_dot_split}
\dot{\Phi}(t) &= -\frac{1}{2}\int_0^\infty\sum_{i,k=1}^N\mu_{ik}(s)\big(\dot{\boldsymbol{x}}_i(t)-\dot{\boldsymbol{x}}_k(t)\big)^T\big(\boldsymbol{\eta}_i^t(s)-\boldsymbol{\eta}_k^t(s)\big)\,ds \notag \\
&\quad - \frac{1}{2}\int_0^\infty\sum_{i,k=1}^N\mu_{ik}(s)\big(\boldsymbol{x}_i(t)-\boldsymbol{x}_k(t)\big)^T\big[\big(\boldsymbol{x}_i(t)-\boldsymbol{x}_k(t)\big) - \partial_s\big(\boldsymbol{\eta}_i^t(s)-\boldsymbol{\eta}_k^t(s)\big)\big]\,ds.
\end{align}

Exchanging indices $i \leftrightarrow k$ and applying the definition
$$
\dot{\boldsymbol{x}}_i(t) = \sum_{k=1}^N \int_0^\infty \mu_{ik}(s)(\boldsymbol{\eta}_k^t(s)-\boldsymbol{\eta}_i^t(s))\,ds,
$$
the first term expands to:
{\small
\begin{align}\label{eq:Phi_term1}
&-\frac{1}{2}\int_0^\infty\sum_{i,k=1}^N \mu_{ik}(s)\big(\dot{\boldsymbol{x}}_i(t)-\dot{\boldsymbol{x}}_k(t)\big)^T\big(\boldsymbol{\eta}_i^t(s)-\boldsymbol{\eta}_k^t(s)\big)\,ds \notag \\
&= \sum_{i=1}^N \dot{\boldsymbol{x}}_i(t)^T \left( -\sum_{k=1}^N \int_0^\infty \mu_{ik}(s)\big(\boldsymbol{\eta}_i^t(s)-\boldsymbol{\eta}_k^t(s)\big)\,ds \right) = \sum_{i=1}^N \dot{\boldsymbol{x}}_i(t)^T \dot{\boldsymbol{x}}_i(t) = \sum_{i=1}^N \|\dot{\boldsymbol{x}}_i(t)\|^2.
\end{align}
}

For the second term in \eqref{eq:Phi_dot_split}, splitting the integrand and applying integration by parts on $\partial_s(\boldsymbol{\eta}_i^t(s)-\boldsymbol{\eta}_k^t(s))$:
{\small
\begin{align}\label{eq:Phi_term2}
&-\frac{1}{2}\sum_{i,k=1}^N \left(\int_0^\infty \mu_{ik}(s)\,ds\right) \|\boldsymbol{x}_i(t)-\boldsymbol{x}_k(t)\|^2 \notag \\
&\quad + \frac{1}{2}\int_0^\infty \sum_{i,k=1}^N \mu_{ik}(s) (\boldsymbol{x}_i(t)-\boldsymbol{x}_k(t))^T \partial_s (\boldsymbol{\eta}_i^t(s)-\boldsymbol{\eta}_k^t(s))\,ds \notag \\
&= -\frac{1}{2}\sum_{i,k=1}^N \lambda_{ik} \|\boldsymbol{x}_i(t)-\boldsymbol{x}_k(t)\|^2 + \frac{1}{2}\sum_{i,k=1}^N (\boldsymbol{x}_i(t)-\boldsymbol{x}_k(t))^T \left[ \mu_{ik}(s)(\boldsymbol{\eta}_i^t(s)-\boldsymbol{\eta}_k^t(s)) \right]_0^\infty \notag \\
&\quad - \frac{1}{2}\int_0^\infty \sum_{i,k=1}^N \mu'_{ik}(s) (\boldsymbol{x}_i(t)-\boldsymbol{x}_k(t))^T (\boldsymbol{\eta}_i^t(s)-\boldsymbol{\eta}_k^t(s))\,ds \notag \\
&= -\frac{1}{2}\sum_{i,k=1}^N \lambda_{ik} \|\boldsymbol{x}_i(t)-\boldsymbol{x}_k(t)\|^2 - \frac{1}{2}\int_0^\infty \sum_{i,k=1}^N \mu'_{ik}(s) (\boldsymbol{x}_i(t)-\boldsymbol{x}_k(t))^T (\boldsymbol{\eta}_i^t(s)-\boldsymbol{\eta}_k^t(s))\,ds,
\end{align}
}
where $\lambda_{ik} := \int_0^\infty \mu_{ik}(s)\,ds = g_{ik}(0) > 0$. Combining \eqref{eq:Phi_term1} and \eqref{eq:Phi_term2}:
{\small
\begin{equation*}\label{eq:Phi_dot_final}
\dot{\Phi}(t) = \sum_{i=1}^N \|\dot{\boldsymbol{x}}_i(t)\|^2 - \frac{1}{2}\sum_{i,k=1}^N \lambda_{ik}\|\boldsymbol{x}_i(t)-\boldsymbol{x}_k(t)\|^2 - \frac{1}{2}\int_0^\infty \sum_{i,k=1}^N \mu'_{ik}(s)(\boldsymbol{x}_i(t)-\boldsymbol{x}_k(t))^T(\boldsymbol{\eta}_i^t(s)-\boldsymbol{\eta}_k^t(s))\,ds.
\end{equation*}
}

Define the total Lyapunov functional $\mathcal{E}(t) := E_0(t) + \varepsilon \Phi(t)$ for $\varepsilon > 0$. Since $\mu'_{ik}(s) \le 0$, writing $-\mu'_{ik}(s) = |\mu'_{ik}(s)| \ge 0$, we apply Young's inequality with parameter $\alpha > 0$:
{\small
\begin{align}\label{eq:cross_bound}
&\frac{\varepsilon}{2}\int_0^\infty \sum_{i,k=1}^N (-\mu'_{ik}(s))(\boldsymbol{x}_i(t)-\boldsymbol{x}_k(t))^T(\boldsymbol{\eta}_i^t(s)-\boldsymbol{\eta}_k^t(s))\,ds \notag \\
&\le \frac{\varepsilon}{4\alpha}\int_0^\infty \sum_{i,k=1}^N (-\mu'_{ik}(s)) \|\boldsymbol{x}_i(t)-\boldsymbol{x}_k(t)\|^2\,ds + \frac{\varepsilon\alpha}{4} \int_0^\infty \sum_{i,k=1}^N (-\mu'_{ik}(s)) \|\boldsymbol{\eta}_i^t(s)-\boldsymbol{\eta}_k^t(s)\|^2\,ds \notag \\
&= \frac{\varepsilon}{4\alpha} \sum_{i,k=1}^N \mu_{ik}(0) \|\boldsymbol{x}_i(t)-\boldsymbol{x}_k(t)\|^2 + \frac{\varepsilon\alpha}{4} \int_0^\infty \sum_{i,k=1}^N (-\mu'_{ik}(s)) \|\boldsymbol{\eta}_i^t(s)-\boldsymbol{\eta}_k^t(s)\|^2\,ds,
\end{align}
}
where we used $\int_0^\infty (-\mu'_{ik}(s))\,ds = \mu_{ik}(0) - \lim\limits_{t\rightarrow\infty}\mu_{ik}(t) = \mu_{ik}(0)$.

Next, applying the Cauchy--Schwarz inequality to $\dot{\boldsymbol{x}}_i(t) = \sum\limits_{k=1}^N \int_0^\infty \mu_{ik}(s)(\boldsymbol{\eta}_k^t(s)-\boldsymbol{\eta}_i^t(s))\,ds$:
\begin{align}\label{eq:xdot_bound}
\|\dot{\boldsymbol{x}}_i(t)\|^2 &= \left\| \sum_{k=1}^N \int_0^\infty \frac{\mu_{ik}(s)}{\sqrt{-\mu'_{ik}(s)}} \cdot \sqrt{-\mu'_{ik}(s)}\, (\boldsymbol{\eta}_k^t(s)-\boldsymbol{\eta}_i^t(s))\,ds \right\|^2 \notag \\
&\le \left(\sum_{k=1}^N \int_0^\infty \frac{\mu_{ik}(s)^2}{-\mu'_{ik}(s)}\,ds\right) \left(\sum_{k=1}^N \int_0^\infty (-\mu'_{ik}(s))\|\boldsymbol{\eta}_i^t(s)-\boldsymbol{\eta}_k^t(s)\|^2\,ds\right).
\end{align}
By Assumption \ref{assum:g}, $\mu_{ik}(s) \le \frac{1}{\delta}(-\mu'_{ik}(s))$, so the first parenthesis is bounded by:
\begin{equation*}
\sum_{k=1}^N \int_0^\infty \frac{\mu_{ik}(s)^2}{-\mu'_{ik}(s)}\,ds \le \frac{1}{\delta}\sum_{k=1}^N \int_0^\infty \mu_{ik}(s)\,ds = \frac{1}{\delta}\sum_{k=1}^N \lambda_{ik} \le \frac{N\lambda_{\max}}{\delta} =: \Lambda < \infty.
\end{equation*}
Summing \eqref{eq:xdot_bound} over $i \in [N]$ yields:
\begin{equation}\label{eq:sum_xdot_bound}
\varepsilon \sum_{i=1}^N \|\dot{\boldsymbol{x}}_i(t)\|^2 \le \varepsilon \Lambda \int_0^\infty \sum_{i,k=1}^N (-\mu'_{ik}(s)) \|\boldsymbol{\eta}_i^t(s)-\boldsymbol{\eta}_k^t(s)\|^2\,ds.
\end{equation}

Substituting \eqref{eq:cross_bound} and \eqref{eq:sum_xdot_bound} into $\dot{\mathcal{E}}(t) = \dot{E}_0(t) + \varepsilon \dot{\Phi}(t)$:
\begin{align}\label{eq:E_dot_assembled}
\dot{\mathcal{E}}(t) &\le -\varepsilon \sum_{i,k=1}^N \left( \frac{\lambda_{ik}}{2} - \frac{\mu_{ik}(0)}{4\alpha} \right) \|\boldsymbol{x}_i(t)-\boldsymbol{x}_k(t)\|^2 \notag \\
&\quad - \int_0^\infty \sum_{i,k=1}^N (-\mu'_{ik}(s)) \|\boldsymbol{\eta}_i^t(s)-\boldsymbol{\eta}_k^t(s)\|^2 \left( \frac{1}{4} - \frac{\varepsilon\alpha}{4} - \varepsilon\Lambda \right)\,ds.
\end{align}

Setting $\lambda_{\min} := \min\limits_{i,k}\lambda_{ik} > 0$ and $\mu^0_{\max} := \max_{i,k}\mu_{ik}(0) < \infty$, we choose $\alpha > 0$ sufficiently large such that
\begin{equation*}
\frac{\lambda_{ik}}{2} - \frac{\mu_{ik}(0)}{4\alpha} \ge \frac{\lambda_{\min}}{2} - \frac{\mu^0_{\max}}{4\alpha} =: c_x > 0, \qquad \forall i,k \in [N].
\end{equation*}
With $\alpha$ fixed, we choose $\varepsilon > 0$ sufficiently small such that
\begin{equation*}
c_\eta := \frac{1}{4} - \varepsilon\left( \frac{\alpha}{4} + \Lambda \right) > 0.
\end{equation*}

Finally, we establish topological equivalence between $\mathcal{E}(t)$ and $E_0(t)$. Applying Cauchy--Schwarz to $\Phi(t)$:
\begin{flalign*}\label{eq:Phi_equiv}
|\Phi(t)| &\le \frac{1}{2}\int_0^\infty \sum_{i,k=1}^N \mu_{ik}(s) \left( \frac{1}{2}\|\boldsymbol{x}_i(t)-\boldsymbol{x}_k(t)\|^2 + \frac{1}{2}\|\boldsymbol{\eta}_i^t(s)-\boldsymbol{\eta}_k^t(s)\|^2 \right) ds \notag \\
&= \frac{1}{4}\sum_{i,k=1}^N \lambda_{ik} \|\boldsymbol{x}_i(t)-\boldsymbol{x}_k(t)\|^2 + \frac{1}{4}\int_0^\infty \sum_{i,k=1}^N \mu_{ik}(s) \|\boldsymbol{\eta}_i^t(s)-\boldsymbol{\eta}_k^t(s)\|^2 ds \notag \\
&\le \lambda_{\max} \left( \frac{N}{2}\sum_{i=1}^N \|\boldsymbol{x}_i(t)\|^2 \right) + \frac{1}{4}\int_0^\infty \sum_{i,k=1}^N \mu_{ik}(s) \|\boldsymbol{\eta}_i^t(s)-\boldsymbol{\eta}_k^t(s)\|^2 ds \notag \\
&\le \max\{N\lambda_{\max}, 1\} E_0(t) =: C_\Phi E_0(t).
\end{flalign*}
By restricting $\varepsilon \le \frac{1}{2 C_\Phi}$, we ensure $|\varepsilon \Phi(t)| \le \frac{1}{2} E_0(t)$, establishing the uniform bounds:
\begin{equation*}\label{eq:equivalence0}
\frac{1}{2} E_0(t) \le \mathcal{E}(t) \le \frac{3}{2} E_0(t), \qquad \forall t \ge 0.
\end{equation*}

Using $-\mu'_{ik}(s) \ge \delta \mu_{ik}(s)$ from Remark \ref{rem:kernel_implication}, \eqref{eq:E_dot_assembled} yields
\begin{align*}
\dot{\mathcal{E}}(t) &\le -2Nc_x \sum_{i=1}^N \|\boldsymbol{x}_i(t)\|^2 - 4\delta c_\eta \left( \frac{1}{4}\int_0^\infty \sum_{i,k=1}^N \mu_{ik}(s)\|\boldsymbol{\eta}_i^t(s)-\boldsymbol{\eta}_k^t(s)\|^2\,ds \right) \notag \\
&\le -\min\{4Nc_x, 4\delta c_\eta\} E_0(t) \le -\frac{2}{3}\min\{4Nc_x, 4\delta c_\eta\} \mathcal{E}(t) =: -\gamma \mathcal{E}(t).
\end{align*}
By Gr\"onwall's Lemma, $\mathcal{E}(t) \le \mathcal{E}(0)e^{-\gamma t}$, and hence $E_0(t) \le 2\mathcal{E}(t) \le 3E_0(0)e^{-\gamma t}$. This establishes unconditional global exponential consensus for Model I.



\section{Exponential flocking for Second-Order Systems}\label{sec:proof_2nd}

In this section, we establish the exponential flocking behaviors for the second-order Cucker--Smale type models, providing the detailed proofs for Theorem \ref{thm:model4} and Theorem \ref{thm:model3}. The primary analytical difficulty for second-order systems lies in the potential expansion of the spatial diameter, which weakens the instantaneous alignment forces. For Model IV, we address this issue by introducing a novel Lyapunov position-integral functional that explicitly couples the velocity energy with the running maximum of the spatial diameter. Utilizing Dini derivatives, we establish a continuous bootstrap framework that simultaneously guarantees uniform spatial boundedness and exponential velocity alignment under a sharp capacity condition. For Model III, where spatial coupling is entirely absent, we construct a velocity-based multiplier functional to extract dissipation purely from the fading memory, proving unconditional exponential flocking.

\subsection{Exponential flocking for system \eqref{eq:model4}}
By Lemma \ref{lem:uncoupled_transform_2nd}, Model IV can be rewritten as
{\small
\begin{equation}\label{eq:model4_auto}
\begin{cases}
\dot{\boldsymbol{x}}_i(t) = \boldsymbol{v}_i(t), \\[4pt]
\dot{\boldsymbol{v}}_i(t) = \displaystyle\sum_{k=1}^{N} \psi(\|\boldsymbol{x}_i(t)-\boldsymbol{x}_k(t)\|)\bigl(\boldsymbol{v}_k(t)-\boldsymbol{v}_i(t)\bigr) + \sum_{k=1}^{N} \int_0^\infty \mu_{ik}(s)\bigl(\boldsymbol{\xi}_k^t(s)-\boldsymbol{\xi}_i^t(s)\bigr)\,ds.
\end{cases}
\end{equation}
}
%Summing the velocity equation over $i \in [N]$ gives $\frac{d}{dt}\sum_{i=1}^N \boldsymbol{v}_i(t) = \mathbf{0}$. In the center-of-mass frame, we set $\sum_{i=1}^N \boldsymbol{v}_i(0) = \mathbf{0}$, whence $\sum_{i=1}^N \boldsymbol{v}_i(t) \equiv \mathbf{0}$ for all $t \ge 0$.
We define the velocity energy functional $E_v(t)$:
\begin{equation}\label{eq:Ev}
E_v(t) := \frac12\sum_{i=1}^N\|\boldsymbol{v}_i(t)\|^2 + \frac14\int_0^\infty\sum_{i,k=1}^N \mu_{ik}(s)\|\boldsymbol{\xi}_i^t(s)-\boldsymbol{\xi}_k^t(s)\|^2\,ds.
\end{equation}
We next recall the definition of the upper right and lower right Dini derivative of a continuous function $f$ at a point $t$:
\begin{equation*}\label{dini}
D^+f(t)=\underset{\Delta t>0,\Delta t\rightarrow0}{\limsup}\frac{f\left(t+\Delta t\right)-f(t)}{\Delta t},~~D_+f(t)=\underset{\Delta t>0,\Delta t\rightarrow0}{\liminf}\frac{f\left(t+\Delta t\right)-f(t)}{\Delta t}.
\end{equation*}
\begin{lemma}{\rm (see \cite[Lemma 15.17]{dini-2})}\label{dini-1}
Given differentiable functions $f_1,...,f_m:(a,b)\rightarrow\R,$ the max function $f_{\max}(t)=\max\left\{f_i(t)|i\in[m]\right\}$ satisfies
$$
D^+f_{\max}(t)=\max\left\{\frac{d}{dt}f_i(t)\bigg|i\in {\rm argmax}\left(f_{\max}(t)\right)\right\}.
$$
\end{lemma}

%\begin{remark}\label{r-dini}
%{\rm From $-D_+f(t)=D^+(-f(t)),$ we can define the lower right Dini derivative of the min function $f_{\min}(t)=\min\left\{f_i(t)|i\in[m]\right\}$ from Lemma \ref{dini-1} as follows:
%\begin{equation*}
%D_+f_{\min}(t)=\min\left\{\frac{d}{dt}f_i(t)\bigg|i\in {\rm argmin}\left(f_{\min}(t)\right)\right\}.
%\end{equation*}}
%\end{remark}


\noindent {\bf Proof of Theorem \ref{thm:model4}.} The proof relies on the construction of a non-increasing Lyapunov position-integral functional. We detail the proof in several consecutive steps to rigorously address the variations in the spatial diameter.

\noindent\textit{Step 1: Energy Dissipation Inequality.}
We first evaluate the time derivative of the velocity energy functional $E_v(t)$ defined in \eqref{eq:Ev}. Differentiating the state velocity kinetic component along the trajectories of \eqref{eq:model4_auto} and symmetrically exchanging the indices $i \leftrightarrow k$ yields:
\begin{align}\label{eq:Ev_dot_state}
\frac{d}{dt}\left(\frac12\sum_{i=1}^N\|\boldsymbol{v}_i(t)\|^2\right) &= \sum_{i=1}^N \boldsymbol{v}_i(t)^T \dot{\boldsymbol{v}}_i(t) \notag \\
&= \sum_{i=1}^N \boldsymbol{v}_i(t)^T \Biggl( \sum_{k=1}^{N} \psi(\|\boldsymbol{x}_i(t)-\boldsymbol{x}_k(t)\|)\bigl(\boldsymbol{v}_k(t)-\boldsymbol{v}_i(t)\bigr) \notag \\
&\qquad\qquad\qquad + \sum_{k=1}^{N} \int_0^\infty \mu_{ik}(s)\bigl(\boldsymbol{\xi}_k^t(s)-\boldsymbol{\xi}_i^t(s)\bigr)\,ds \Biggr) \notag \\
&= -\frac12\sum_{i,k=1}^N\psi(\|\boldsymbol{x}_i(t)-\boldsymbol{x}_k(t)\|)\|\boldsymbol{v}_i(t)-\boldsymbol{v}_k(t)\|^2 \notag\\
&\quad -\frac12\int_0^\infty\sum_{i,k=1}^N\mu_{ik}(s)\bigl(\boldsymbol{v}_i(t)-\boldsymbol{v}_k(t)\bigr)^T\bigl(\boldsymbol{\xi}_i^t(s)-\boldsymbol{\xi}_k^t(s)\bigr)\,ds.
\end{align}

Differentiating the velocity memory energy utilizing the transport identity $\partial_t\boldsymbol{\xi}_i^t(s) = \boldsymbol{v}_i(t) - \partial_s\boldsymbol{\xi}_i^t(s)$, we obtain:
\begin{align}\label{eq:Ev_dot_memory}
&\frac{d}{dt}\left(\frac14\int_0^\infty\sum_{i,k=1}^N\mu_{ik}(s)\|\boldsymbol{\xi}_i^t(s)-\boldsymbol{\xi}_k^t(s)\|^2\,ds\right) \notag \\
&= \frac12\int_0^\infty\sum_{i,k=1}^N\mu_{ik}(s)\bigl(\boldsymbol{\xi}_i^t(s)-\boldsymbol{\xi}_k^t(s)\bigr)^T \bigl( \partial_t\boldsymbol{\xi}_i^t(s) - \partial_t\boldsymbol{\xi}_k^t(s) \bigr)\,ds \notag \\
&= \frac12\int_0^\infty\sum_{i,k=1}^N\mu_{ik}(s)\bigl(\boldsymbol{\xi}_i^t(s)-\boldsymbol{\xi}_k^t(s)\bigr)^T \left[ \bigl(\boldsymbol{v}_i(t)-\boldsymbol{v}_k(t)\bigr) - \partial_s\bigl(\boldsymbol{\xi}_i^t(s)-\boldsymbol{\xi}_k^t(s)\bigr) \right] ds \notag \\
&= \frac12\int_0^\infty\sum_{i,k=1}^N\mu_{ik}(s)\bigl(\boldsymbol{v}_i(t)-\boldsymbol{v}_k(t)\bigr)^T\bigl(\boldsymbol{\xi}_i^t(s)-\boldsymbol{\xi}_k^t(s)\bigr)\,ds \notag \\
&\quad - \frac14\int_0^\infty\sum_{i,k=1}^N \mu_{ik}(s)\,\partial_s\|\boldsymbol{\xi}_i^t(s)-\boldsymbol{\xi}_k^t(s)\|^2\,ds.
\end{align}

Summing \eqref{eq:Ev_dot_state} and \eqref{eq:Ev_dot_memory}, and applying integration by parts to the last term of \eqref{eq:Ev_dot_memory}, we arrive at:
{\small
\begin{align}\label{eq:Ev_dot_final}
\dot{E}_v(t) &= -\frac12\sum_{i,k=1}^N\psi(\|\boldsymbol{x}_i(t)-\boldsymbol{x}_k(t)\|)\|\boldsymbol{v}_i(t)-\boldsymbol{v}_k(t)\|^2 \notag \\
&\quad - \frac{1}{4}\sum_{i,k=1}^N \left[ \mu_{ik}(s) \|\boldsymbol{\xi}_i^t(s) - \boldsymbol{\xi}_k^t(s)\|^2 \right]_0^\infty + \frac14\int_0^\infty\sum_{i,k=1}^N \mu_{ik}'(s)\|\boldsymbol{\xi}_i^t(s)-\boldsymbol{\xi}_k^t(s)\|^2\,ds \notag \\
&= -\frac12\sum_{i,k=1}^N\psi(\|\boldsymbol{x}_i(t)-\boldsymbol{x}_k(t)\|)\|\boldsymbol{v}_i(t)-\boldsymbol{v}_k(t)\|^2 + \frac14\int_0^\infty\sum_{i,k=1}^N \mu_{ik}'(s)\|\boldsymbol{\xi}_i^t(s)-\boldsymbol{\xi}_k^t(s)\|^2\,ds\nonumber\\
&\le -2N\psi(X_{\max}(t))\left(\frac12\sum_{i=1}^N\|\boldsymbol{v}_i(t)\|^2\right) - \delta\left(\frac14\int_0^\infty\sum_{i,k=1}^N \mu_{ik}(s)\|\boldsymbol{\xi}_i^t(s)-\boldsymbol{\xi}_k^t(s)\|^2\,ds\right) \notag \\
&\le -\min\bigl\{2N\psi(X_{\max}(t)),\, \delta\bigr\} E_v(t),
\end{align}
}
where in the second to last inequality, we used the decay assumption $\mu_{ik}'(s)\le-\delta\mu_{ik}(s)$, the non-increasing property of $\psi$ which implies $\psi(\|\boldsymbol{x}_i(t)-\boldsymbol{x}_k(t)\|) \ge \psi(X_{\max}(t))$, and the algebraic identity $\sum\limits_{i,k=1}^N \|\boldsymbol{v}_i(t)-\boldsymbol{v}_k(t)\|^2 = 2N \sum\limits_{i=1}^N \|\boldsymbol{v}_i(t)\|^2$.

\noindent\textit{Step 2: Evolution of the Square Root of Energy.}
If $E_v(t_0) = 0$ at some time $t_0 \ge 0$, then in the center-of-mass frame we have $\boldsymbol{v}_i(t_0) = \mathbf{0}$ and historical relative displacements $\boldsymbol{\xi}_i^{t_0}(s)-\boldsymbol{\xi}_k^{t_0}(s) = \mathbf{0}$ for all $i,k \in [N], s \ge 0$. By the uniqueness of solutions to \eqref{eq:model4_auto}, the particles permanently freeze into a static state with $\boldsymbol{v}_i(t) \equiv \mathbf{0}$ and $X_{\max}(t) \equiv X_{\max}(t_0) < \infty$ for all $t \ge t_0$, for which the asymptotic flocking holds trivially.

Thus, without loss of generality, we assume $E_v(t) > 0$ for all $t \ge 0$. Differentiating $\sqrt{E_v(t)}$ and using \eqref{eq:Ev_dot_final}, we obtain:
\begin{equation}\label{eq:sqrt_Ev_dot}
\frac{d}{dt}\sqrt{E_v(t)} = \frac{\dot{E}_v(t)}{2\sqrt{E_v(t)}} \le -\frac{1}{2} \min\bigl\{2N\psi(X_{\max}(t)),\, \delta\bigr\} \sqrt{E_v(t)}, \qquad \forall t \ge 0.
\end{equation}

\noindent\textit{Step 3: Dini Derivative of the Running Maximum.}
From the equation $\dot{\boldsymbol{x}}_i(t) = \boldsymbol{v}_i(t)$, applying Lemma \ref{dini-1} to the spatial diameter $X_{\max}(t) = \max\limits_{i,k\in[N]} \|\boldsymbol{x}_i(t) - \boldsymbol{x}_k(t)\|$ yields:
\begin{flalign}\label{eq:Dini_Xmax}
 D^+ X_{\max}(t) & = \max_{(i,k)\in\operatorname{argmax}(X_{\max}(t))} \frac{(\boldsymbol{x}_i(t) - \boldsymbol{x}_k(t))^T(\boldsymbol{v}_i(t) - \boldsymbol{v}_k(t))}{\|\boldsymbol{x}_i(t) - \boldsymbol{x}_k(t)\|}\nonumber\\
 &\leq \max_{(i,k)\in\operatorname{argmax}(X_{\max}(t))} \|\bv_i(t)-\bv_k(t)\|\nonumber\le 2 \max_{j\in[N]} \|\boldsymbol{v}_j(t)\|\leq 2\sqrt{\sum_{i\in[N]}\|\bv_i(t)\|^2}\nonumber\\
   &  \le 2\sqrt{2}\sqrt{E_v(t)}.
\end{flalign}
We define the running maximum function $Y(t)$ by:
\begin{equation*}\label{eq:Y_def}
Y(t) := \max_{s \in [0,t]} X_{\max}(s), \quad t \ge 0.
\end{equation*}
Clearly, $Y(t)$ is continuous and monotonically non-decreasing, which guarantees for all $t \ge 0,$
$$
D^+ Y(t) \ge 0.
$$
Moreover, for any $t \ge 0$ and $\Delta t > 0$, we have
$$
Y(t+\Delta t) = \max\left\{ Y(t), \max_{s\in[t, t+\Delta t]} X_{\max}(s) \right\}.
$$
If $X_{\max}(t) < Y(t)$, continuity implies $Y(t+\Delta t) = Y(t)$ for small $\Delta t > 0$, whence $D^+ Y(t) = 0$.
If $X_{\max}(t) = Y(t)$, from \eqref{eq:Dini_Xmax} we have
$$
D^+ Y(t)= D^+X_{\max}(t)\le V_{\max}(t) \le 2\sqrt{2}\sqrt{E_v(t)}.
$$
Therefore, for all $t \ge 0,$
\begin{equation}\label{eq:Y_Dini_bound}
0 \le D^+ Y(t) \le 2\sqrt{2}\sqrt{E_v(t)}.
\end{equation}
Since $X_{\max}(t) \le Y(t)$ and $\psi$ is non-increasing, $\psi(X_{\max}(t)) \ge \psi(Y(t))$. Substituting this and \eqref{eq:Y_Dini_bound} into \eqref{eq:sqrt_Ev_dot}:
\begin{flalign}\label{eq:sqrt_Ev_dot_Y}
 \frac{d}{dt}\sqrt{E_v(t)}  & \le -\frac{1}{2} \min\bigl\{2N\psi(Y(t)),\, \delta\bigr\} \sqrt{E_v(t)}\nonumber\\
  &\leq-\frac{1}{4\sqrt{2}} \min\bigl\{2N\psi(Y(t)),\, \delta\bigr\} D^+ Y(t).
\end{flalign}
\noindent\textit{Step 4: Construction of the Lyapunov Functional.}
We define the Lyapunov position-integral functional $L(t)$ by:
\begin{equation*}\label{eq:L_func_def}
L(t) := \sqrt{E_v(t)} + \frac{N}{2\sqrt{2}} \int_{X_{\max}(0)}^{Y(t)} \min\left\{\psi(r),\, \frac{\delta}{2N}\right\} dr, \quad t \ge 0.
\end{equation*}
Applying the chain rule for Dini derivatives to the $C^1$-smooth integral term, we obtain from \eqref{eq:sqrt_Ev_dot_Y}:
\begin{equation*}\label{eq:Dini_L_nonpositive}
D^+ L(t) = \frac{d}{dt}\sqrt{E_v(t)} + \frac{N}{2\sqrt{2}} \min\left\{\psi(Y(t)),\, \frac{\delta}{2N}\right\} D^+ Y(t) \le 0, \quad \forall t \ge 0.
\end{equation*}
Due to the continuity, $L(t)$ is monotonically non-increasing. Hence $L(t) \le L(0) = \sqrt{E_v(0)}$.
Since $\sqrt{E_v(t)} \ge 0$, this yields the uniform bound for all $t \ge 0,$
\begin{equation*}\label{eq:integral_bound}
\frac{N}{2\sqrt{2}} \int_{X_{\max}(0)}^{Y(t)} \min\left\{\psi(r),\, \frac{\delta}{2N}\right\} dr \le \sqrt{E_v(0)}.
\end{equation*}

\noindent\textit{Step 5: Uniform Boundedness of Spatial Diameter.}
Define the spatial capacity function for $y \ge X_{\max}(0),$
$$
\Psi(y) := \frac{N}{2\sqrt{2}} \int_{X_{\max}(0)}^{y} \min\left\{\psi(r),\, \frac{\delta}{2N}\right\} dr.
$$
Since the integrand is strictly positive, $\Psi(y)$ is strictly monotonically increasing.
By the capacity condition \eqref{eq:flocking_capacity_cond},
$$
\lim_{y \to \infty} \Psi(y) > \sqrt{E_v(0)}.
$$
Because $\Psi(Y(t)) \le \sqrt{E_v(0)}$ for all $t \ge 0$, there exists a unique finite root $R_* > X_{\max}(0)$ to the algebraic equation
\begin{equation*}\label{eq:R_star_def}
\frac{N}{2\sqrt{2}} \int_{X_{\max}(0)}^{R_*} \min \left\{ \psi(r),\, \frac{\delta}{2N} \right\} dr = \sqrt{E_v(0)}.
\end{equation*}
Hence
\begin{equation*}\label{eq:Xmax_global_bound}
\sup_{t \ge 0} X_{\max}(t) \le \sup_{t \ge 0} Y(t) \le R_* < \infty.
\end{equation*}

\noindent\textit{Step 6: Exponential Velocity Alignment.}
Since $X_{\max}(t) \le R_*$ for all $t \ge 0$, we have $\psi(X_{\max}(t)) \ge \psi(R_*) > 0$. Inequality \eqref{eq:Ev_dot_final} simplifies to:
\begin{equation*}\label{eq:Ev_exponential}
\dot{E}_v(t) \le -\gamma_v(R_*) E_v(t), \quad \text{where } \gamma_v(R_*) := \min\bigl\{2N\psi(R_*),\, \delta\bigr\} > 0.
\end{equation*}
By Gr\"onwall's Lemma, $E_v(t) \le E_v(0) e^{-\gamma_v(R_*) t}$ for all $t \ge 0$, which immediately implies for all $t \ge 0,$
\begin{equation*}
V_{\max}(t) \le 2\max_{i\in[N]}\|\boldsymbol{v}_i(t)\| \le 2\sqrt{2E_v(t)} \le 2\sqrt{2E_v(0)}\, e^{-\frac{1}{2}\gamma_v(R_*) t}.
\end{equation*}
This completes the proof of Theorem \ref{thm:model4}.


\subsection{Exponential flocking for system \eqref{eq:model3}}
For Model III ($\psi \equiv 0$), equation \eqref{eq:Ev_dot_final} reduces to
$$
\dot{E}_v(t) = \frac14\int_0^\infty\sum_{i,k=1}^N \mu_{ik}'(s)\|\boldsymbol{\xi}_i^t(s)-\boldsymbol{\xi}_k^t(s)\|^2\,ds \le 0.
$$
We introduce the velocity multiplier functional:
\begin{equation}\label{eq:Phi_v}
\Phi_v(t) := -\frac12\int_0^\infty\sum_{i,k=1}^N \mu_{ik}(s)\bigl(\boldsymbol{v}_i(t)-\boldsymbol{v}_k(t)\bigr)^T\bigl(\boldsymbol{\xi}_i^t(s)-\boldsymbol{\xi}_k^t(s)\bigr)\,ds.
\end{equation}

\noindent {\bf Proof of Theorem \ref{thm:model3}.} Differentiating \eqref{eq:Phi_v} along \eqref{eq:model4_auto} with $\psi \equiv 0$ using $\partial_t \boldsymbol{\xi}_i^t(s) = \boldsymbol{v}_i(t) - \partial_s \boldsymbol{\xi}_i^t(s)$:
\begin{align}\label{eq:Phiv_dot_split}
\dot{\Phi}_v(t) &= -\frac12\int_0^\infty\sum_{i,k=1}^N\mu_{ik}(s)\bigl(\dot{\boldsymbol{v}}_i(t)-\dot{\boldsymbol{v}}_k(t)\bigr)^T\bigl(\boldsymbol{\xi}_i^t(s)-\boldsymbol{\xi}_k^t(s)\bigr)\,ds \notag \\
&\quad -\frac12\int_0^\infty\sum_{i,k=1}^N\mu_{ik}(s)\bigl(\boldsymbol{v}_i(t)-\boldsymbol{v}_k(t)\bigr)^T\Bigl[\bigl(\boldsymbol{v}_i(t)-\boldsymbol{v}_k(t)\bigr)-\partial_s\bigl(\boldsymbol{\xi}_i^t(s)-\boldsymbol{\xi}_k^t(s)\bigr)\Bigr]\,ds.
\end{align}

Exchanging indices $i \leftrightarrow k$ and substituting $\dot{\boldsymbol{v}}_i(t) = \sum\limits_{k=1}^N \int_0^\infty \mu_{ik}(s)(\boldsymbol{\xi}_k^t(s)-\boldsymbol{\xi}_i^t(s))\,ds$, the first integral evaluates to:
{\small
\begin{align}\label{eq:Phiv_term1}
&-\frac12\int_0^\infty \sum_{i,k=1}^N \mu_{ik}(s) \big(\dot{\boldsymbol{v}}_i(t)-\dot{\boldsymbol{v}}_k(t)\big)^T \big(\boldsymbol{\xi}_i^t(s)-\boldsymbol{\xi}_k^t(s)\big)\,ds \notag \\
&= \sum_{i=1}^N \dot{\boldsymbol{v}}_i(t)^T \left( -\sum_{k=1}^N \int_0^\infty \mu_{ik}(s)(\boldsymbol{\xi}_i^t(s)-\boldsymbol{\xi}_k^t(s))\,ds \right) = \sum_{i=1}^N \dot{\boldsymbol{v}}_i(t)^T \dot{\boldsymbol{v}}_i(t) = \sum_{i=1}^N \|\dot{\boldsymbol{v}}_i(t)\|^2.
\end{align}
}

For the second integral in \eqref{eq:Phiv_dot_split}, splitting the terms and applying integration by parts on $\partial_s(\boldsymbol{\xi}_i^t(s)-\boldsymbol{\xi}_k^t(s))$, we obtain:
{\footnotesize
\begin{equation}\label{eq:Phiv_term2}
\begin{array}{l}
\displaystyle{-\frac12\int_0^\infty\sum_{i,k=1}^N\mu_{ik}(s)\bigl(\boldsymbol{v}_i(t)-\boldsymbol{v}_k(t)\bigr)^T\Bigl[\bigl(\boldsymbol{v}_i(t)-\boldsymbol{v}_k(t)\bigr)-\partial_s\bigl(\boldsymbol{\xi}_i^t(s)-\boldsymbol{\xi}_k^t(s)\bigr)\Bigr]\,ds\hspace{1cm}}\\
\displaystyle{\hspace{1cm}=-\frac12\sum_{i,k=1}^N \Big(\int_0^\infty \mu_{ik}(s)\,ds\Big) \|\boldsymbol{v}_i(t)-\boldsymbol{v}_k(t)\|^2} \\
\displaystyle{\hspace{3cm}+ \frac12\int_0^\infty \sum_{i,k=1}^N \mu_{ik}(s)(\boldsymbol{v}_i(t)-\boldsymbol{v}_k(t))^T \partial_s (\boldsymbol{\xi}_i^t(s)-\boldsymbol{\xi}_k^t(s))\,ds  }\\
%&= -\frac12\sum_{i,k=1}^N \lambda_{ik}\|\boldsymbol{v}_i(t)-\boldsymbol{v}_k(t)\|^2 + \frac12\sum_{i,k=1}^N (\boldsymbol{v}_i(t)-\boldsymbol{v}_k(t))^T \left[ \mu_{ik}(s)(\boldsymbol{\xi}_i^t(s)-\boldsymbol{\xi}_k^t(s)) \right]_0^\infty \notag \\
%&\quad - \frac12\int_0^\infty\sum_{i,k=1}^N\mu_{ik}'(s)\bigl(\boldsymbol{v}_i(t)-\boldsymbol{v}_k(t)\bigr)^T\bigl(\boldsymbol{\xi}_i^t(s)-\boldsymbol{\xi}_k^t(s)\bigr)\,ds \notag \\
\displaystyle{\hspace{1cm}= -\frac12\sum_{i,k=1}^N \lambda_{ik}\|\boldsymbol{v}_i(t)-\boldsymbol{v}_k(t)\|^2 - \frac12\int_0^\infty\sum_{i,k=1}^N\mu_{ik}'(s)\bigl(\boldsymbol{v}_i(t)-\boldsymbol{v}_k(t)\bigr)^T\bigl(\boldsymbol{\xi}_i^t(s)-\boldsymbol{\xi}_k^t(s)\bigr)\,ds,}
\end{array}
\end{equation}
}
where $\lambda_{ik} = \int_0^\infty \mu_{ik}(s)\,ds = g_{ik}(0) > 0$. Combining \eqref{eq:Phiv_term1} and \eqref{eq:Phiv_term2}:
{\small
\begin{equation*}\label{eq:Phiv_dot_final}
\dot{\Phi}_v(t) = \sum_{i=1}^N\|\dot{\boldsymbol{v}}_i(t)\|^2 -\frac12\sum_{i,k=1}^N\lambda_{ik}\|\boldsymbol{v}_i(t)-\boldsymbol{v}_k(t)\|^2 -\frac12\int_0^\infty\sum_{i,k=1}^N\mu_{ik}'(s)\bigl(\boldsymbol{v}_i(t)-\boldsymbol{v}_k(t)\bigr)^T\bigl(\boldsymbol{\xi}_i^t(s)-\boldsymbol{\xi}_k^t(s)\bigr)\,ds.
\end{equation*}
}


Define the total velocity Lyapunov functional $\mathcal{E}_v(t) := E_v(t) + \varepsilon \Phi_v(t)$ for $\varepsilon > 0$. Applying Young's inequality with parameter $\alpha > 0$:
\begin{align}\label{eq:v_cross_bound}
&\frac{\varepsilon}{2}\int_0^\infty\sum_{i,k=1}^N\bigl(-\mu_{ik}'(s)\bigr)\bigl(\boldsymbol{v}_i(t)-\boldsymbol{v}_k(t)\bigr)^T\bigl(\boldsymbol{\xi}_i^t(s)-\boldsymbol{\xi}_k^t(s)\bigr)\,ds \notag \\
&\le \frac{\varepsilon}{4\alpha}\sum_{i,k=1}^N\mu_{ik}(0)\|\boldsymbol{v}_i(t)-\boldsymbol{v}_k(t)\|^2 +\frac{\varepsilon\alpha}{4}\int_0^\infty\sum_{i,k=1}^N\bigl(-\mu_{ik}'(s)\bigr)\|\boldsymbol{\xi}_i^t(s)-\boldsymbol{\xi}_k^t(s)\|^2\,ds.
\end{align}

Applying Cauchy--Schwarz to $\dot{\boldsymbol{v}}_i(t) = \sum\limits_{k=1}^N \int_0^\infty \mu_{ik}(s)(\boldsymbol{\xi}_k^t(s)-\boldsymbol{\xi}_i^t(s))\,ds$:
\begin{align}\label{eq:vdot_CS}
\|\dot{\boldsymbol{v}}_i(t)\|^2 &= \left\| \sum_{k=1}^N \int_0^\infty \frac{\mu_{ik}(s)}{\sqrt{-\mu'_{ik}(s)}} \cdot \sqrt{-\mu'_{ik}(s)}\, (\boldsymbol{\xi}_k^t(s)-\boldsymbol{\xi}_i^t(s))\,ds \right\|^2 \notag \\
&\le \left(\sum_{k=1}^N\int_0^\infty\frac{\mu_{ik}(s)^2}{-\mu_{ik}'(s)}\,ds\right)\left(\sum_{k=1}^N\int_0^\infty(-\mu_{ik}'(s))\|\boldsymbol{\xi}_i^t(s)-\boldsymbol{\xi}_k^t(s)\|^2\,ds\right).
\end{align}
Since $\mu_{ik}(s) \le \frac{1}{\delta}(-\mu'_{ik}(s))$, the first parenthesis is bounded by
$$
\frac{1}{\delta}\sum_{k=1}^N \lambda_{ik} \le \frac{N\lambda_{\max}}{\delta} =: K^v < \infty.
$$
Summing \eqref{eq:vdot_CS} over $i \in [N]$:
\begin{equation}\label{eq:sum_vdot_bound}
\varepsilon\sum_{i=1}^N\|\dot{\boldsymbol{v}}_i(t)\|^2 \le \varepsilon K^v\int_0^\infty\sum_{i,k=1}^N(-\mu_{ik}'(s))\|\boldsymbol{\xi}_i^t(s)-\boldsymbol{\xi}_k^t(s)\|^2\,ds.
\end{equation}

Combining \eqref{eq:v_cross_bound} and \eqref{eq:sum_vdot_bound} into $\dot{\mathcal{E}}_v(t) = \dot{E}_v(t) + \varepsilon \dot{\Phi}_v(t)$:
\begin{flalign*}\label{eq:Ev_dot_assembled}
\dot{\mathcal{E}}_v(t) &\le -\varepsilon\sum_{i,k=1}^N\left(\frac{\lambda_{ik}}{2}-\frac{\mu_{ik}(0)}{4\alpha}\right)\|\boldsymbol{v}_i(t)-\boldsymbol{v}_k(t)\|^2 \notag\\
&\quad -\int_0^\infty\sum_{i,k=1}^N(-\mu_{ik}'(s))\|\boldsymbol{\xi}_i^t(s)-\boldsymbol{\xi}_k^t(s)\|^2 \left(\frac14-\frac{\varepsilon\alpha}{4}-\varepsilon K^v\right)ds.
\end{flalign*}

Choosing $\alpha > 0$ sufficiently large such that $\frac{\lambda_{ik}}{2}-\frac{\mu_{ik}(0)}{4\alpha} \ge c_v > 0$, and choosing $\varepsilon > 0$ small enough such that $c_\xi := \frac14-\frac{\varepsilon\alpha}{4}-\varepsilon K^v > 0$ and $|\varepsilon\Phi_v(t)| \le \frac12 E_v(t)$, we obtain the equivalence $\frac12 E_v(t) \le \mathcal{E}_v(t) \le \frac32 E_v(t)$ and
\begin{equation*}
\dot{\mathcal{E}}_v(t) \le -\gamma_v \mathcal{E}_v(t)
\end{equation*}
for a uniform rate $\gamma_v > 0$.

By Gr\"onwall's Lemma, $E_v(t) \le 3 E_v(0) e^{-\gamma_v t}$ unconditionally for all $t \ge 0$. Integrating $\dot{\boldsymbol{x}}_i(t) = \boldsymbol{v}_i(t)$ gives:
\begin{align*}
\|\boldsymbol{x}_i(t)-\boldsymbol{x}_k(t)\| &\le \|\boldsymbol{x}_i(0)-\boldsymbol{x}_k(0)\| + \int_0^t \|\boldsymbol{v}_i(\tau)-\boldsymbol{v}_k(\tau)\|\,d\tau \\
&\le X_{\max}(0) + 2\sqrt{2}\int_0^t \sqrt{E_v(\tau)}\,d\tau \\
&\le X_{\max}(0) + 2\sqrt{6E_v(0)}\int_0^t e^{-\frac{\gamma_v}{2} \tau}\,d\tau \\
&\le X_{\max}(0) + \frac{4\sqrt{6E_v(0)}}{\gamma_v} < \infty, \qquad \forall t \ge 0.
\end{align*}
Thus $\sup\limits_{t\ge0} X_{\max}(t) < \infty$, completing the proof of Theorem \ref{thm:model3}.

\section{Conclusion}\label{sec:conclusion}

In this paper, we analyzed the emergent collective behaviors of multi-agent systems driven by an infinite distributed fading memory of Volterra type. By introducing exact Dafermos past-history transformations, the original integro-differential models were rigorously reformulated into dynamical systems defined on extended product Hilbert spaces. For the first-order opinion dynamics, we proved that fading memory intrinsically induces a hidden dissipative mechanism. By constructing auxiliary multiplier functionals to extract this dissipation, we established that global exponential consensus occurs unconditionally, even in the complete absence of instantaneous communication forces. For the second-order velocity alignment model, we provided both unconditional and conditional flocking frameworks. In the pure fading memory regime, we showed that exponential velocity alignment and spatial boundedness emerge unconditionally. When instantaneous communication is present, we established an explicit integral capacity condition connecting the initial phase configurations, the memory decay rate, and the communication weight. Using suitable Lyapunov functionals, we proved the exponential emergence of flocking.

There are several remaining issues that we did not investigate in this work. For example, our current analysis fundamentally relies on symmetric interaction topologies and uniform fading memory kernels among all agents. In reality, the interaction topology is often given by a general digraph, and different agents may possess heterogeneous cognitive memory rates or be subjected to communication noise. Thus, it would be interesting and challenging to extend our results to scenarios with asymmetric couplings, directed network topologies, or heterogeneous time-delays. Rigorous analysis for these remaining issues will be addressed in future research.


\begin{thebibliography}{99}

%\bibitem{Acebron2005}
%J.~A.~Acebr\'on, L.~L.~Bonilla, C.~J.~P\'erez~Vicente, F.~Ritort, and R.~Spigler, ``The Kuramoto model: A simple paradigm for synchronization phenomena,'' \textit{Rev. Mod. Phys.}, 77 (2005), 137--185.


\bibitem{Blondel2009}
V.~D.~Blondel, J.~M.~Hendrickx, and J.~N.~Tsitsiklis, ``On Krause's multi-agent consensus model with state-dependent connectivity'', \textit{IEEE Trans. Automat. Control}, 54 (2009), 2586--2597.

\bibitem{Boschi2021}
G.~Boschi, C.~Cammarota, and R.~K\"uhn, ``Opinion dynamics with emergent collective memory: the impact of a long and heterogeneous news history'', \textit{Phys. A}, 569 (2021), Paper No. 125799, 19 pp.



%\bibitem{Buck1966}
%J.~Buck and E.~Buck, ``Biology of synchronous flashing of fireflies,'' \textit{Nature}, 211 (1966), 562--564.

\bibitem{dini-2}
F. Bullo,
\newblock \emph{Lectures on Network Systems},
\newblock 1.6$^{nd}$ edition, Kindle Direct Publishing, 2022.

\bibitem{Bullo2009}
F.~Bullo, J.~Cort\'es, and S.~Mart\'inez, \textit{Distributed Control of Robotic Networks: A Mathematical Approach to Motion Coordination Algorithms}, Princeton University Press, Princeton, NJ, 2009.


\bibitem{Camazine2001}
S.~Camazine, J.~L.~Deneubourg, N.~R.~Franks, J.~Sneyd, G.~Theraulaz, and E.~Bonabeau, \textit{Self-Organization in Biological Systems}, Princeton University Press, Princeton, NJ, 2001.

\bibitem{Canuto2012}
C.~Canuto, F.~Fagnani, and P.~Tilli, ``An Eulerian approach to the analysis of Krause's consensus models'', \textit{SIAM J. Control Optim.}, 50 (2012), 243--265.

\bibitem{Carrillo2010}
J.~A.~Carrillo, M.~Fornasier, G.~Toscani, and F.~Vecil, ``Particle, kinetic, and hydrodynamic models of swarming'', in \textit{Mathematical Modeling of Collective Behavior in Socio-Economic and Life Sciences}, Birkh\"auser, Boston, 2010, pp. 297--336.

\bibitem{Castellano2009}
C.~Castellano, S.~Fortunato, and V.~Loreto, ``Statistical physics of social dynamics'', \textit{Rev. Mod. Phys.}, 81 (2009), 591--646.

\bibitem{Chepyzhov2006}
V.~V.~Chepyzhov, E.~Mainini, and V.~Pata, ``Stability of abstract linear semigroups arising from heat conduction with memory'', \textit{Asymptot. Anal.}, 50 (2006), 269--291.

\bibitem{Cho2025}
H.~Cho, S.-Y.~Ha, and M.~Kang, ``Emergent behaviors of a Kuramoto ensemble under fading memory'', \textit{J. Nonlinear Sci.}, 35 (2025), Art. 9.

\bibitem{Chiara}
Y.-P.~Choi, C.~Cicolani, and C.~Pignotti,
``Time-delayed opinion dynamics with leader-follower interactions: consensus, stability, and mean-field limits'',
Commun. Math. Sci. 24 (2026), no. 4, 1073--1102.

\bibitem{Choi2021}
Y.-P.~Choi, A.~Paolucci, and C.~Pignotti, ``Consensus of the Hegselmann--Krause opinion formation model with time delay'', \textit{Math. Methods Appl. Sci.}, 44 (2021), 4560--4579.

\bibitem{Conti2020}
M.~Conti, F.~Dell'Oro, and V.~Pata, ``Exponential decay of a first order linear Volterra equation'', \textit{Math. Eng.}, 2 (2020), 459--471.

\bibitem{Continelli}
E.~Continelli,
``Asymptotic flocking for the Cucker-Smale model with time variable time delays'',
\textit{Acta Appl. Math.}, 188 (2023), Paper No. 15, 23 pp.

\bibitem{Continelli2023}
E.~Continelli and C.~Pignotti, ``Consensus for Hegselmann--Krause type models with time variable time delays'', \textit{Math. Methods Appl. Sci.}, 46 (2023), 18916--18934.

\bibitem{Cucker2007}
F.~Cucker and S.~Smale, ``Emergent behavior in flocks'', \textit{IEEE Trans. Automat. Control}, 52 (2007), 852--862.

\bibitem{Dafermos1970}
C.~M.~Dafermos, ``Asymptotic stability in viscoelasticity'', \textit{Arch. Ration. Mech. Anal.}, 37 (1970), 297--308.

%\bibitem{Dong2013}
%J.-G.~Dong and X.~Xue, ``Synchronization %analysis of Kuramoto oscillators,'' %\textit{Commun. Math. Sci.}, 11 (2013), 465--480.
\bibitem{Dong}
J.-G.~Dong, S.-Y.~Ha, J.~Jung, Jinwook, and D.~Kim,
 ``On the stochastic flocking of the Cucker-Smale flock with randomly switching topologies'',
\textit{SIAM J. Control Optim.} 58 (2020), 2332--2353.

\bibitem{Du2025}
L.~Du, S.-Y.~Ha, and H.~Yu, ``Interplay of heterogeneous time-delays and feedback control in the Hegselmann--Krause model'', \textit{J. Differential Equations}, 423 (2025), 597--630.

\bibitem{Hask_memory}
R. Erban and J.~Haskovec, ``Impact of memory on clustering in spontaneous particle aggregation'',  arXiv preprint arXiv:2510.15335 (2025).

\bibitem{Giorgi2001}
C.~Giorgi, J.~E.~Mu\~{n}oz Rivera, and V.~Pata, ``Global attractors for a semilinear hyperbolic equation in viscoelasticity'', \textit{J. Math. Anal. Appl.}, 260 (2001), 83--99.


\bibitem{Ha2009}
S.-Y.~Ha and J.-G.~Liu, ``A simple proof of the  Cucker-Smale flocking
dynamics and mean-field limit'', \textit{Commun. Math. Sci.}, 2 (2009), 297--325.


\bibitem{Ha2}
S.-Y. Ha and E. Tadmor,
``From particle to kinetic and hydrodynamic descriptions of
	flocking'',
\textit{Kinet. Relat. Models}, 1, (2008), 415--435.


%\bibitem{Ha2020}
%S.-Y.~Ha and S.-Y.~Ryoo, ``Asymptotic %phase-locking dynamics and critical coupling %strength for the Kuramoto model,'' %\textit{Commun. Math. Phys.}, 377 (2020), %811--857.

\bibitem{Haskovec2021}
J.~Haskovec, ``A simple proof of asymptotic consensus in the Hegselmann--Krause and Cucker--Smale models with normalization and delay'', \textit{SIAM J. Appl. Dyn. Syst.}, 20 (2021), 130--148.



\bibitem{Hegselmann2002}
R.~Hegselmann and U.~Krause, ``Opinion dynamics and bounded confidence: models, analysis and simulation'', \textit{J. Artif. Soc. Soc. Simul.}, 5 (2002), no. 3, Art. 2.


\bibitem{Jabin2014}
P.-E.~Jabin and S.~Motsch, ``Clustering and asymptotic behavior in opinion formation'', \textit{J. Differential Equations}, 257 (2014), 4165--4187.

\bibitem{Jackson2008}
M.~O.~Jackson, \textit{Social and Economic Networks}, Princeton University Press, Princeton, NJ, 2008.

\bibitem{Jadbabaie2003}
A.~Jadbabaie, J.~Lin, and A.~S.~Morse, ``Coordination of groups of mobile autonomous agents using nearest neighbor rules'', \textit{IEEE Trans. Automat. Control}, 48 (2003), 988--1001.

\bibitem{Jiang2025}
M.~Jiang, W.~Su, G.~Ren, and Y.~Yu, ``Memory-driven bounded confidence opinion dynamics: A Hegselmann--Krause model based on fractional-order methods'', arXiv preprint arXiv:2506.04701, (2025).

%\bibitem{Kuramoto1975}
%Y.~Kuramoto, ``Self-entrainment of a %population of coupled non-linear %oscillators,'' \textit{Lect. Notes %Phys.}, 39 (1975), 420--422.

\bibitem{Liu2023}
Q.~Liu and L.~Chai, ``The memory influence on opinion dynamics in coopetitive social networks: Analysis, application, and simulation'', \textit{IEEE Trans. Control Netw. Syst.}, 10 (2023), 1867--1878.

\bibitem{Liu2026}
Q. Liu, A. Li, and L. Chai, ``Dynamics of opinion propagation with memory and fake news'', \textit{Automatica J. IFAC}, 185 (2026), Paper No. 112745, 14 pp.

%\bibitem{Lorenz2007}
%J.~Lorenz, ``Continuous opinion dynamics under bounded confidence: A survey,'' \textit{Int. J. Mod. Phys. C}, 18 (2007), 1819--1838.
%\bibitem{Nicaise2006}
%S.~Nicaise and C.~Pignotti, ``Stabilization of the wave equation with variable coefficients and boundary condition of memory type,'' \textit{Asymptot. Anal.}, 50 (2006), 31--67.

\bibitem{OlfatiSaber2007}
R.~Olfati-Saber, J.~A.~Fax, and R.~M.~Murray, ``Consensus and cooperation in networked multi-agent systems'', \textit{Proc. IEEE}, 95 (2007), 215--233.

\bibitem{Pignotti2018}
C.~Pignotti and E.~Tr\'elat, ``Convergence to consensus of the general finite-dimensional Cucker--Smale model with time-varying delays'', \textit{Commun. Math. Sci.}, 16 (2018), 2053--2076.

%\bibitem{Pikovsky2001}
%A.~Pikovsky, M.~Rosenblum, and J.~Kurths, \textit{Synchronization: A Universal Concept in Nonlinear Sciences}, Cambridge University Press, Cambridge, 2001.

\bibitem{Cartabia}
M. Rodriguez Cartabia,
``Cucker-Smale model with time delay'',
\newblock{\em Discrete Contin. Dynam. Systems}, 42 (2022), no. 5, 2409--2432.

\bibitem{Su2017}
W.~Su, G.~Chen, and Y.~Hong, ``Noise leads to quasi-consensus of Hegselmann--Krause opinion dynamics'', \textit{Automatica}, 85 (2017), 448--454.

\bibitem{Vicsek1995}
T.~Vicsek, A.~Czir\'ok, E.~Ben-Jacob, I.~Cohen, and O.~Shochet, ``Novel type of phase transition in a system of self-driven particles'', \textit{Phys. Rev. Lett.}, 75 (1995), 1226--1229.

\bibitem{Volterra1930}
%V.~Volterra, \textit{Th\'eorie G\'en\'erale des Functionnelles et des \'Equations Int\'egrales et Integro-Diff\'erentielles}, Gauthier-Villars, Paris, 1930.
V. Volterra, \emph{Theory of functionals and of integral and integro-differential equations}, Dover, New York, 1959.


\bibitem{wang2025agent}
Z. Z. Wang, J. Mao, D. Fried, and G. Neubig, ``Agent Workflow Memory'', in  \textit{Forty-second International Conference on Machine Learning}, (2025), https://openreview.net/forum?id=NTAhi2JEEE.



%\bibitem{Winfree1967}
%A.~T.~Winfree, ``Biological rhythms and the behavior of populations of coupled oscillators,'' \textit{J. Theor. Biol.}, 16 (1967), 15--42.


\bibitem{Fanqin}
F.~Zeng, X.~Xue, and Y.~Zhu,
``Critical exponent for Cucker-Smale model under group-hierarchical multi-leadership'',
\textit{Appl. Math. Lett.}, 136 (2023),  No. 108452, 7 pp.

\end{thebibliography}

\end{document} 
